% concatenated from 4_Final_Cohomological_Reduction.tex

\documentclass[11pt,a4paper]{article}

\usepackage[T1]{fontenc}
\usepackage[utf8]{inputenc}
\usepackage{lmodern}
\usepackage{amsmath,amssymb,amsthm,mathtools}
\usepackage{geometry}
\usepackage{microtype}
\usepackage{enumitem}
\usepackage{hyperref}
\usepackage{mathrsfs}
\usepackage{bm}
\geometry{margin=1in}
\hypersetup{colorlinks=true,linkcolor=blue,citecolor=blue,urlcolor=blue}

\newtheorem{theorem}{Theorem}[section]
\newtheorem{lemma}[theorem]{Lemma}
\newtheorem{proposition}[theorem]{Proposition}
\newtheorem{corollary}[theorem]{Corollary}
\newtheorem{remark}[theorem]{Remark}
\newtheorem{definition}[theorem]{Definition}
\newtheorem{assumption}[theorem]{Assumption}

\newcommand{\R}{\mathbb R}
\newcommand{\M}{\mathcal M}
\newcommand{\Lip}{\operatorname{Lip}}
\newcommand{\diam}{\operatorname{diam}}
\newcommand{\htop}{h_{\mathrm{top}}}
\newcommand{\Per}{\operatorname{Per}}
\newcommand{\dist}{\operatorname{dist}}
\newcommand{\supp}{\operatorname{supp}}
\newcommand{\varphihat}{\widehat{\varphi}}
\newcommand{\Sfn}{S_n}
\newcommand{\eps}{\varepsilon}

% Author affiliation commands (compatible with the article class).
\newcommand{\address}[1]{\def\theaddress{#1}}
\newcommand{\email}[1]{\def\theemail{#1}}
\newcommand{\printauthorinfo}{%
  \begin{center}
  \small
  \theaddress\\
  \href{mailto:\theemail}{\texttt{\theemail}}
  \end{center}
}

\title{Cohomological Reduction for Fiber-Contracting Extensions:\\
From Subcohomology to Thermodynamic Formalism}
\author{G. Ferreira}
\address{G. Ferreira, Departamento de Matem\'atica, Universidade Federal do Maranh\~ao, Av. dos Portugueses 1966, Vila Bacanga, 65065--545 S\~ao Lu\'is, MA, Brazil}
\email{giovane.ferreira@ufma.br}
\date{}

\begin{document}
\maketitle
\printauthorinfo

\begin{center}
\small
\textbf{2020 Mathematics Subject Classification.}
37D35, 37A60, 37D20, 37C30.

\textbf{Keywords.}
cohomological equation; subcohomology; Liv\v{s}ic theorem;
fiber contraction; ergodic optimization; thermodynamic formalism;
weak Gibbs measures; sequential Gibbs measures.
\end{center}

\begin{abstract}
We develop a reduction principle for cohomological, variational, and
thermodynamic questions in fiber-contracting extensions of local
homeomorphisms. Under uniform contraction in the fibers and the existence
of a continuous global section, every H\"older potential admits an explicit
cohomological reduction
\[
\varphi=\psi\circ\pi+u-u\circ F,
\]
where the reduced potential $\psi$ depends only on the quotient. The
construction is quantitative and separates the hypotheses required for
continuity, H\"older regularity, subcohomology, Liv\v{s}ic theory, weak
Gibbs transfer, and sequential Gibbs transfer.

The main contribution developed in this paper is a transfer framework: the quotient theory can be
lifted to the attractor while preserving invariant averages, minimizing
measures, metric entropy, pressure, and equilibrium states. In particular,
the invariant-measure spaces are in bijection and the contracting fibers
have zero relative topological entropy. We also establish weak Gibbs
transfer under a precise subexponential comparison of Bowen-ball masses and
give concrete sufficient criteria. We also prove blockwise Bowen sequential weak Gibbs transfer
along point-dependent Gibbs times; under uniformly bounded sequential mass
distortion, the strong Bowen sequential Gibbs property is preserved along the same prescribed admissible Gibbs-time sequence.

The paper emphasizes explicit applications. We work out an affine
contracting skew-product
\[
F(x,y)=(f(x),ay+\rho(x)),\qquad |a|<1,
\]
for which fiber-dependent potentials of the form
$\varphi(x,y)=g(x)+by$ are reduced explicitly to
\[
\psi(x)=g(x)+\frac{b}{1-a}\rho(x).
\]
This example shows concretely how a potential that genuinely depends on the
stable coordinate can be replaced by a quotient potential and how
equilibrium states and ergodic optimization are consequently reduced to the
base dynamics. A second symbolic skew-product example verifies the weak
Gibbs correspondence on an invariant graph. These examples demonstrate
that the framework is not merely formal and clarify precisely where the
additional geometric hypotheses enter.
\end{abstract}

\section{Introduction}

Cohomological equations and subcohomological inequalities are among the
basic tools for identifying the dynamically relevant part of an observable.
For a dynamical system $T:X\to X$, the equation
\[
\phi=u-u\circ T
\]
implies that Birkhoff sums telescope and that the integral of $\phi$ against
every invariant probability measure vanishes. Liv\v{s}ic theory identifies
periodic orbit obstructions as the decisive obstruction in a broad
hyperbolic setting, while subactions and calibrated subactions turn
variational inequalities into pointwise inequalities
\cite{Livsic,deLaLlaveMarcoMoriyon,BouschJenkinson,MbarkiSantana}.

The purpose of this paper is to establish a systematic reduction mechanism
for these questions in a class of fiber-contracting extensions. Let
\[
F:N\to N,\qquad f:M\to M
\]
be continuous maps and let $\pi:N\to M$ be a continuous surjection satisfying
\[
\pi\circ F=f\circ\pi.
\]
For $x\in M$, write $N_x=\pi^{-1}(x)$ and assume that points in the same
fiber converge exponentially under $F$. We also assume that $\pi$ admits a
continuous global section. The latter is a genuine topological hypothesis;
it is not automatic for an arbitrary factor map and will be verified
explicitly in the principal examples below.

The central theorem says that the contracting direction is
\emph{cohomologically removable}. For every H\"older potential $\varphi$,
the stable series
\[
u(\hat x)=\sum_{j=0}^{\infty}
\left[\varphi(F^j\hat x)-\varphi(F^jr(\hat x))\right]
\]
converges uniformly and gives
\[
\boxed{\varphi=\psi\circ\pi+u-u\circ F.}
\]
Thus the potential may be replaced, within its cohomology class, by one
that is constant on the fibers. This mechanism is closely related to the
Sinai cohomology lemma in hyperbolic symbolic models
\cite{Sinai}; the point here is not to claim the stable reduction itself as
a new version of Sinai's lemma, but to formulate it for abstract
fiber-contracting extensions and to derive a unified set of consequences.

The main structural message is therefore
\[
\boxed{
\begin{array}{c}
\text{fiber contraction}\\
\Downarrow\\
\text{cohomological reduction}\\
\Downarrow\\
\text{transfer of quotient theory to the attractor}
\end{array}}
\]
Once the reduction is available, the invariant-measure problem is completely
encoded by the quotient: projection gives a bijection between invariant
probability measures, invariant averages agree, and minimizing measures are
carried bijectively to minimizing measures. Consequently, ergodic
optimization for a fiber-dependent potential is reduced to an optimization
problem for the quotient potential.

The same reduction has thermodynamic consequences. Uniform contraction
implies zero fiberwise relative topological entropy. We use the
Ledrappier--Walters relative variational principle in the form
\[
\sup_{\pi_*\widetilde\nu=\mu} h_{\widetilde\nu}(F)
=
h_\mu(f)+\int_M h_{\mathrm{top}}(F,N_x)\,d\mu(x).
\]
Since metric entropy is monotone under factors, every invariant lift also
satisfies $h_\mu(f)\le h_{\widetilde\mu}(F)$. Once the fiber entropies are
shown to vanish, the supremum above equals $h_\mu(f)$; hence every invariant
lift has exactly the same metric entropy as its projection. Combining this equality with the
cohomological identity gives
\[
P_F(\varphi)=P_f(\psi).
\]
Equilibrium states therefore correspond exactly under the invariant-measure
bijection. The fibers contribute no additional exponential orbit complexity,
even though the original potential may depend strongly on the fiber
coordinate. We stress that this conclusion uses the relative variational
principle in addition to the reduction theorem; it is not a formal
consequence of cohomology alone.

A second part of the paper concerns subcohomology and Liv\v{s}ic theory.
The critical case $m(\varphi,F)=0$ is transferred from a corresponding
subcohomological theorem on the quotient, whereas the positive-margin case
$m(\varphi,F)>0$ admits a direct Ma\~n\'e-type construction. Periodic points
of the quotient lift uniquely to periodic points of the extension by the
contraction of the return map on each periodic fiber. Hence quotient
Liv\v{s}ic theorems lift without introducing a new periodic obstruction.
For zooming quotients this yields, in particular, a direct route from the
subcohomology and Liv\v{s}ic theory of Mbarki--Santana to the fibered
attractor \cite{MbarkiSantana}.

Weak Gibbs transfer requires more than zero relative entropy. We therefore
isolate the precise additional condition: for the particular corresponding
invariant measures under consideration, Bowen-ball masses must be comparable
up to a subexponential factor, possibly after changing the Bowen radius.
This condition is not implied by fiber contraction alone: contraction controls
relative orbit separation inside a single fiber, whereas a Bowen ball also
contains information about how mass is distributed across nearby fibers.
Accordingly, the weak Gibbs results are stated as conditional transfer
theorems, not as automatic consequences of the reduction. We prove forward
and backward transfer theorems and give concrete sufficient criteria in
invariant-graph, contracting skew-product, and local-product settings.

The concrete examples are a central part of the paper. First, we analyze
the affine skew-product
\[
F(x,y)=\bigl(f(x),ay+\rho(x)\bigr),\qquad |a|<1,
\]
on a compact cylinder. For the genuinely fiber-dependent potential
\[
\varphi(x,y)=g(x)+by
\]
the transfer function and reduced potential can be computed explicitly:
\[
u(x,y)=\frac{b}{1-a}y,
\qquad
\psi(x)=g(x)+\frac{b}{1-a}\rho(x).
\]
Thus the stable coordinate does not disappear by assuming the potential is
stable-constant; rather, its contribution is converted into an explicit
quotient term. This example provides a direct test of the main theorem and
makes the pressure, equilibrium-state, and ergodic-optimization reductions
fully concrete.

Second, we consider a two-sided symbolic skew-product with an invariant
graph. In that setting the invariant lift is supported on the graph, so
Bowen-ball comparison has no exponential loss. Consequently any Gibbs
measure for the quotient potential lifts to a weak Gibbs measure for the
original fiber-dependent potential. This verifies the weak Gibbs mechanism
in a genuinely dynamical model rather than leaving it as an abstract
measure-comparison hypothesis.

These examples also clarify the scope of the theory. Smale--Williams
solenoids and partially hyperbolic attractors provide important motivation,
but the global-section assumption is not automatic in those settings.
Accordingly, we distinguish verified examples from conditional classes and
do not claim an application unless all hypotheses are checked.

For clarity, we stress the precise novelty claim. The stable cohomological
reduction is a Sinai-type mechanism and is not claimed as new by itself. The
contribution developed in the present paper is the abstract transfer framework
for fiber-contracting extensions, together with its simultaneous consequences
for invariant optimization, subcohomology, Liv\v{s}ic theory, pressure,
equilibrium states, weak Gibbs measures, and sequential Gibbs transfer, as well
as the explicit zooming-extension application in
Section~\ref{subsec:zooming-application}.

The paper is organized around this division of labor. Sections
\ref{sec:geom}--\ref{sec:reduction} establish the geometric and
cohomological reduction; the subsequent sections transfer subcohomology,
Liv\v{s}ic theory, invariant optimization, and thermodynamic formalism.
Section~\ref{sec:weak-gibbs} isolates the additional quantitative measure
hypothesis, and Section~\ref{sec:sequential-gibbs} treats sequential Gibbs
transfer along point-dependent Gibbs times. The final section contains the verified examples and explains
how the abstract results apply to concrete fibered models.

\section{Geometric Setting}\label{sec:geom}

Let $N$ and $M$ be compact metric spaces. Let $F:N\to N$ be continuous and
$f:M\to M$ be a local homeomorphism. Assume that $\pi:N\to M$ is a
continuous surjection satisfying
\begin{equation}
    \pi\circ F=f\circ\pi.
\end{equation}
Then
\begin{equation}
    F(N_x)\subset N_{f(x)}.
    \label{eq:fiber-map}
\end{equation}

\begin{assumption}[Uniform fiber contraction]
There exists $0<\lambda<1$ such that
\[
d_N(F^n\hat x,F^n\hat y)
\leq\lambda^n d_N(\hat x,\hat y)
\]
for all $n\geq0$ and all $\hat x,\hat y$ in the same fiber.
\end{assumption}

\begin{assumption}[Continuous fiber selector]
\label{ass:selector}
There exists a continuous map
\[
r:N\to N
\]
such that
\begin{equation}
r(\hat x)\in N_{\pi(\hat x)}
\end{equation}
and
\begin{equation}
\pi(\hat x)=\pi(\hat y)
\quad\Longrightarrow\quad
r(\hat x)=r(\hat y).
\end{equation}
\end{assumption}

Thus $r$ is equivalently a continuous section of $\pi$ written on $N$.
Indeed, since $\pi$ is a continuous surjection from the compact space $N$
onto the Hausdorff space $M$, it is a quotient map. Because $r$ is continuous
and constant on the fibers of $\pi$, it factors uniquely through a continuous
map $s:M\to N$ such that
\[
\pi\circ s=\operatorname{id}_M,
\qquad
r=s\circ\pi.
\]
Conversely, every continuous section $s$ defines such a selector by
$r=s\circ\pi$. This makes clear that the selector assumption is strong and is
not automatic for a general factor map. It must be verified in concrete
examples.

For $0<\beta\leq1$, define
\[
[\varphi]_\beta
=
\sup_{\hat x\neq\hat y}
\frac{|\varphi(\hat x)-\varphi(\hat y)|}
{d_N(\hat x,\hat y)^\beta}.
\]

We shall also use invariant measure spaces
\[
\M_F(N)=\{\widetilde\mu:F_*\widetilde\mu=\widetilde\mu\},
\qquad
\M_f(M)=\{\mu:f_*\mu=\mu\}.
\]
For a continuous potential $\varphi$ define
\begin{equation}
m(\varphi,F)
=
\min_{\widetilde\mu\in\M_F(N)}
\int_N\varphi\,d\widetilde\mu,
\label{eq:m-total}
\end{equation}
and similarly
\begin{equation}
m(\psi,f)
=
\min_{\mu\in\M_f(M)}
\int_M\psi\,d\mu.
\label{eq:m-base}
\end{equation}
The invariant-measure sets are nonempty by the Krylov--Bogolyubov theorem,
compact in the weak-$*$ topology, and the integral of a continuous potential
is weak-$*$ continuous. Hence the minima in
\eqref{eq:m-total}--\eqref{eq:m-base} exist.

\section{Main Cohomological Reduction}\label{sec:reduction}

\begin{theorem}[Cohomological reduction]
\label{thm:reduction}
Assume the geometric setting above. Let $\varphi:N\to\R$ be
$\beta$-H\"older. Define
\begin{equation}
 u(\hat x)=\sum_{j=0}^{\infty}
 \bigl[\varphi(F^j\hat x)-\varphi(F^jr(\hat x))\bigr].
 \label{eq:u-series-main}
\end{equation}
Then the series converges absolutely and uniformly and defines a continuous
function $u:N\to\R$. Moreover,
\begin{equation}
 \|u\|_\infty\le
 \frac{[\varphi]_\beta\diam(N)^\beta}{1-\lambda^\beta}.
 \label{eq:u-sup-bound}
\end{equation}
The function
\[
 \bar\varphi:=\varphi-u+u\circ F
\]
is constant on each fiber. Consequently there exists a unique continuous
$\psi:M\to\R$ such that
\begin{equation}
 \boxed{\varphi=\psi\circ\pi+u-u\circ F.}
 \label{eq:main-cohomology}
\end{equation}
H\"older regularity is not asserted from fiber contraction alone; a
quantitative sufficient criterion is given in
Proposition~\ref{prop:holder-reduction}.
\end{theorem}

\begin{proof}
Since $\hat x$ and $r(\hat x)$ lie in the same fiber,
\[
d_N(F^j\hat x,F^jr(\hat x))
\le \lambda^j d_N(\hat x,r(\hat x)).
\]
Hence
\begin{align*}
|\varphi(F^j\hat x)-\varphi(F^jr(\hat x))|
&\le [\varphi]_\beta
 d_N(F^j\hat x,F^jr(\hat x))^\beta\\
&\le [\varphi]_\beta\lambda^{\beta j}
 d_N(\hat x,r(\hat x))^\beta\\
&\le [\varphi]_\beta\lambda^{\beta j}\diam(N)^\beta.
\end{align*}
The Weierstrass $M$-test gives absolute uniform convergence and
\eqref{eq:u-sup-bound}; continuity follows because the partial sums are
continuous.

It remains to prove that
$\bar\varphi=\varphi-u+u\circ F$ is constant on fibers. Let
$\pi(\hat x)=\pi(\hat y)$. Since $r(\hat x)=r(\hat y)$, subtraction of
the two absolutely convergent series gives
\begin{align}
u(\hat x)-u(\hat y)
&=\sum_{j=0}^{\infty}
\bigl[\varphi(F^j\hat x)-\varphi(F^j\hat y)\bigr].
\label{eq:u-fiber-difference}
\end{align}
The points $F\hat x$ and $F\hat y$ also lie in the same fiber, and hence
\begin{align}
u(F\hat x)-u(F\hat y)
&=\sum_{j=0}^{\infty}
\bigl[\varphi(F^{j+1}\hat x)-\varphi(F^{j+1}\hat y)\bigr]\\
&=\sum_{j=1}^{\infty}
\bigl[\varphi(F^j\hat x)-\varphi(F^j\hat y)\bigr].
\label{eq:u-fiber-difference-shifted}
\end{align}
Both series are absolutely convergent because
\[
|\varphi(F^j\hat x)-\varphi(F^j\hat y)|
\le [\varphi]_\beta\lambda^{\beta j}
 d_N(\hat x,\hat y)^\beta.
\]
Subtracting \eqref{eq:u-fiber-difference} and
\eqref{eq:u-fiber-difference-shifted} therefore yields
\begin{align*}
\bar\varphi(\hat x)-\bar\varphi(\hat y)
&=\varphi(\hat x)-\varphi(\hat y)
 -\bigl(u(\hat x)-u(\hat y)\bigr)
 +\bigl(u(F\hat x)-u(F\hat y)\bigr)\\
&=0.
\end{align*}
Thus $\bar\varphi$ is constant on every fiber. Define
\[
\psi(x)=\bar\varphi(s(x)).
\]
Then $\psi$ is continuous and $\bar\varphi=\psi\circ\pi$. For the
canonical transfer function $u$ fixed by \eqref{eq:u-series-main}, the quotient
potential $\psi$ is unique by surjectivity of $\pi$.
\end{proof}

\begin{theorem}[Reduction and transfer principle]
\label{thm:master-reduction}
Assume the geometric setting and the hypotheses of Theorem~\ref{thm:reduction}.
For every H\"older potential $\varphi$ there are a canonically defined continuous
transfer function $u$, given by the stable-difference series
\eqref{eq:u-series-main}, and a corresponding continuous quotient potential
$\psi$ such that
\[
\varphi=\psi\circ\pi+u-u\circ F.
\]
Moreover:
\begin{enumerate}[label=\textnormal{(\roman*)}]
\item $\pi_*:\M_F(N)\to\M_f(M)$ is a bijection;
\item corresponding invariant measures have the same integral of
$\varphi$ and $\psi$, hence
\[
m(\varphi,F)=m(\psi,f);
\]
\item the minimizing measures correspond bijectively;
\item the relative topological entropy is zero and corresponding invariant
measures have equal metric entropy;
\item
\[
P_F(\varphi)=P_f(\psi),
\]
and equilibrium states correspond bijectively;
\item whenever a subcohomological or Liv\v{s}ic theorem on the quotient applies to
$\psi$ (with all of its regularity and dynamical hypotheses satisfied), the
corresponding conclusion transfers to $\varphi$;
\item weak Gibbs transfer holds whenever the corresponding invariant
measures satisfy the forward/backward subexponential Bowen-mass comparison
proved in Section~\ref{sec:weak-gibbs};
\item blockwise sequential weak Gibbs transfer holds along prescribed point-dependent Gibbs times under the sequential Bowen-mass comparison of
Section~\ref{sec:sequential-gibbs}; under uniformly bounded comparison the
strong Bowen sequential Gibbs property is preserved along the same prescribed admissible Gibbs-time sequence.
\end{enumerate}
The conclusions in (i)--(v) use only the hypotheses explicitly identified
in the corresponding sections; the additional assumptions in (vi)--(viii)
are not consequences of fiber contraction alone.
\end{theorem}

\begin{proof}
The reduction formula is Theorem~\ref{thm:reduction}. Item (i) is
Theorem~\ref{thm:lift-bijection}; item (ii) follows from
Corollary~\ref{cor:average-correspondence}; and item (iii) is
Theorem~\ref{thm:ergodic-optimization}. Item (iv) follows from
Lemma~\ref{lem:relative-zero} and Theorem~\ref{thm:pressure}, while item (v)
is the pressure and equilibrium-state correspondence in
Theorem~\ref{thm:pressure}. Item (vi) follows by pulling the quotient
conclusion back through $\pi$ and adding the coboundary $u-u\circ F$;
the Liv\v{s}ic case is made explicit in Theorem~\ref{thm:livsic}. Item (vii) is the content of the forward and backward weak Gibbs transfer
theorems in Section~\ref{sec:weak-gibbs}, while item (viii) is
Theorems~\ref{thm:sequential-weak-forward} and~\ref{thm:sequential-weak-backward}, together with
Corollary~\ref{cor:sequential-gibbs-transfer}. Thus the master principle is a
synthesis of the results established below, with no additional dynamical
hypothesis hidden in the statement.
\end{proof}

\begin{remark}[Canonical normalization and gauge freedom]
\label{rem:gauge-freedom}
The quotient potential is canonical only after fixing the transfer function
$u$ by the stable-difference series \eqref{eq:u-series-main}. Indeed, if
$w\in C(M)$ and
\[
 u'=u+w\circ\pi,
 \qquad
 \psi'=\psi-w+w\circ f,
\]
then semiconjugacy gives the equally valid decomposition
\[
 \varphi=\psi'\circ\pi+u'-u'\circ F.
\]
Thus arbitrary cohomological decompositions have the natural gauge freedom
$(u,\psi)\mapsto(u+w\circ\pi,\psi-w+w\circ f)$. The phrase
``the quotient potential'' below refers to the canonical representative
associated with \eqref{eq:u-series-main}.
\end{remark}

\begin{remark}[Modular use of the transfer principle]
\label{rem:master-modular}
The reduction theorem is the only input needed to remove the fiber dependence
of the potential. The subsequent conclusions are not all consequences of
cohomological reduction alone. In particular, entropy and pressure require
the relative entropy argument of Section~\ref{sec:pressure}, while
subcohomology and Liv\v{s}ic theory require a theorem on the quotient.
Likewise, weak Gibbs transfer requires the quantitative Bowen-mass
hypotheses of Section~\ref{sec:weak-gibbs}. Thus the theorem is to be read as a modular transfer
statement rather than as a single theorem in which all conclusions follow
from fiber contraction alone.
\end{remark}

\begin{remark}
The significance of Theorem~\ref{thm:master-reduction} is not that each
individual conclusion is new in isolation. Rather, a single geometric
reduction converts a fiber-dependent potential into a quotient potential
and thereby makes several theories simultaneously accessible. In
particular, the theorem can be used as a ``preprocessing'' step before
applying a quotient result that is already known for expanding, zooming,
symbolic, or non-uniformly expanding dynamics.
\end{remark}

\begin{proposition}[Summable-defect H\"older criterion]
\label{prop:holder-reduction}
For $j\ge0$ define the fiber defect
\begin{equation}
\Delta_j(\hat x)
:=\varphi(F^j\hat x)-\varphi(F^jr(\hat x)).
\label{eq:defect-delta}
\end{equation}
Assume that, for some $0<\alpha\le1$, each $\Delta_j$ is
$\alpha$-H\"older and that
\begin{equation}
\sum_{j=0}^{\infty}[\Delta_j]_\alpha<\infty.
\label{eq:defect-holder-sum}
\end{equation}
Then the transfer function $u=\sum_{j\ge0}\Delta_j$ is
$\alpha$-H\"older and
\begin{equation}
[u]_\alpha\le\sum_{j=0}^{\infty}[\Delta_j]_\alpha.
\label{eq:u-holder-bound}
\end{equation}
If, in addition, $F$ is Lipschitz, $\alpha\le\beta$, and the section
$s:M\to N$ is $\sigma$-H\"older for some $0<\sigma\le1$, then $\bar\varphi$ is
$\alpha$-H\"older and $\psi=\bar\varphi\circ s$ is
$\alpha\sigma$-H\"older.
\end{proposition}

\begin{proof}
For $\hat x,\hat y\in N$, absolute convergence and
\eqref{eq:defect-holder-sum} give
\begin{align*}
|u(\hat x)-u(\hat y)|
&\le\sum_{j=0}^{\infty}
|\Delta_j(\hat x)-\Delta_j(\hat y)|\\
&\le\left(\sum_{j=0}^{\infty}[\Delta_j]_\alpha\right)
 d_N(\hat x,\hat y)^\alpha.
\end{align*}
This proves \eqref{eq:u-holder-bound}. If $F$ is Lipschitz, then
$u\circ F$ is $\alpha$-H\"older. Since $\alpha\le\beta$ and $N$ is
compact, every $\beta$-H\"older function is also $\alpha$-H\"older;
therefore
$\bar\varphi=\varphi-u+u\circ F$ is $\alpha$-H\"older. Finally, writing
\[
[s]_\sigma
:=\sup_{x\neq y}\frac{d_N(s(x),s(y))}{d_M(x,y)^\sigma},
\]
we obtain
\[
|\psi(x)-\psi(y)|
=|\bar\varphi(s(x))-\bar\varphi(s(y))|
\le [\bar\varphi]_\alpha [s]_\sigma^\alpha
 d_M(x,y)^{\alpha\sigma}.
\]
\end{proof}

\begin{remark}
Condition \eqref{eq:defect-holder-sum} is an analytic transverse
summability criterion, not a consequence of fiber contraction alone. In a
concrete hyperbolic, expanding, or zooming model it should be verified from
the available holonomy, inverse-branch, or distortion estimates. The
advantage of the formulation is that the exact regularity input needed by
the proof is explicit and no unquantified phrase such as ``the same
estimates imply H\"older regularity'' is used.
\end{remark}

\section{Subcohomology}

\subsection{The quotient property}

\begin{definition}[Subcohomological property]
A dynamical system $(M,f)$ has property $\mathrm{SC}$ for a class
$\mathcal C$ of potentials if, for every $\psi\in\mathcal C$,
\[
m(\psi,f)\geq0
\quad\Longrightarrow\quad
\exists v\in C(M):
\qquad
\psi\geq v-v\circ f.
\]
\end{definition}

\begin{theorem}[Subcohomology transfer]
\label{thm:SC-transfer}
Assume the cohomological reduction
\[
\varphi=\psi\circ\pi+u-u\circ F
\]
and suppose $(M,f)$ has property $\mathrm{SC}$ for the class containing
$\psi$. If
\[
m(\varphi,F)\geq0,
\]
then there exists $U\in C(N)$ such that
\[
\boxed{\varphi\geq U-U\circ F.}
\]
If $u$ and $v\circ\pi$ are H\"older, where
\[
\psi\geq v-v\circ f,
\]
then
\[
U=u+v\circ\pi
\]
is H\"older. In particular, this conclusion holds if $v$ is
$\alpha$-H\"older and $\pi$ is $\tau$-H\"older, in which case
$v\circ\pi$ is $\alpha\tau$-H\"older.
\end{theorem}

\begin{proof}
By Corollary~\ref{cor:average-correspondence},
\[
m(\psi,f)=m(\varphi,F)\geq0.
\]
The property $\mathrm{SC}$ gives $v\in C(M)$ satisfying
\[
\psi\geq v-v\circ f.
\]
Set
\[
U=u+v\circ\pi.
\]
Using $\pi\circ F=f\circ\pi$,
\begin{align}
U-U\circ F
&=
u-u\circ F
+
(v-v\circ f)\circ\pi\\
&\leq
u-u\circ F+\psi\circ\pi\\
&=\varphi.
\end{align}
This proves the result.
\end{proof}

\begin{corollary}[Zooming quotients]
\label{cor:zooming}
Suppose $f$ is a zooming system satisfying the hypotheses of
Mbarki--Santana. If the fibered extension satisfies the reduction
hypotheses and the induced potential $\psi$ belongs to the regularity class
required by the quotient subcohomology theorem, then every H\"older potential
$\varphi$ with $m(\varphi,F)\geq0$ admits a continuous subaction
\[
\varphi\geq U-U\circ F.
\]
When H\"older regularity of the lifted subaction is desired, one must in
addition verify the regularity of the pullback $v\circ\pi$. For example, if
the quotient theorem gives an $\alpha$-H\"older subaction $v$ and the
factor map $\pi$ is $\tau$-H\"older, then $v\circ\pi$ is
$\alpha\tau$-H\"older; together with the regularity of $u$ this gives a
H\"older lifted subaction. In the product skew-products considered below,
$\pi$ is Lipschitz.
\end{corollary}

\begin{remark}[Why the critical case cannot be replaced by the elementary
weighted series]
The weighted construction
\[
v=\sum_{j=0}^{\infty}\theta^j|\phi\circ T^j|
\]
gives
\[
v\circ T-v
=
\left(\frac1\theta-1\right)
\sum_{j=1}^{\infty}\theta^j|\phi\circ T^j|
-|\phi|.
\]
The uniform estimate for the first term is
\[
\left(\frac1\theta-1\right)
\frac{\theta}{1-\theta}\|\phi\|_\infty
=
\|\phi\|_\infty.
\]
Thus this term does \emph{not} tend to zero as $\theta\to1^-$. Hence
this weighted-series estimate cannot establish the critical case. The
regime $m(\phi,T)=0$ requires a genuine subcohomological theorem.
\end{remark}

\subsection{Positive margin and a Mañ\'e-type construction}

For a continuous map $T:X\to X$ of a compact metric space and
$\phi\in C(X)$, write
\[
 m(\phi,T)
 :=\min_{\mu\in\mathcal M_T(X)}\int_X\phi\,d\mu.
\]
The minimum exists by compactness of $\mathcal M_T(X)$ and weak-$*$
continuity of the integral. The positive-margin case admits a direct argument
which is conceptually different from the critical case.

\begin{lemma}[Uniform growth of Birkhoff sums]
\label{lem:uniform-growth}
Let $T:X\to X$ be continuous on compact $X$ and $\phi\in C(X)$. If
\[
m(\phi,T)>0,
\]
then there exist $c>0$ and $N_0\in\mathbb N$ such that
\[
S_n\phi(x)
=
\sum_{j=0}^{n-1}\phi(T^jx)
\geq cn
\]
for every $x\in X$ and every $n\geq N_0$.
\end{lemma}

\begin{proof}
Choose $0<c<m(\phi,T)$. Suppose the conclusion is false. Then there are
$n_k\to\infty$ and $x_k\in X$ such that
\[
S_{n_k}\phi(x_k)<cn_k.
\]
Define the empirical measures
\[
\mu_k
=
\frac1{n_k}
\sum_{j=0}^{n_k-1}\delta_{T^jx_k}.
\]
By compactness of the space of probability measures, after passing to a
subsequence $\mu_k\to\mu$ in the weak-$*$ topology. Standard telescoping
gives, for every $\psi\in C(X)$,
\[
\left|
\int\psi\circ T\,d\mu_k-\int\psi\,d\mu_k
\right|
\leq
\frac{2\|\psi\|_\infty}{n_k}\to0.
\]
Hence $\mu$ is $T$-invariant. On the other hand,
\[
\int\phi\,d\mu
=
\lim_{k\to\infty}
\frac{S_{n_k}\phi(x_k)}{n_k}
\leq c<m(\phi,T),
\]
a contradiction.
\end{proof}

\begin{proposition}[Positive-margin subaction]
\label{prop:positive}
If $m(\phi,T)>0$, then there exists $u\in C(X)$ such that
\[
\phi\geq u-u\circ T.
\]
\end{proposition}

\begin{proof}
By Lemma~\ref{lem:uniform-growth}, choose $c>0$ and $N_0$ such that
$S_n\phi(x)\ge cn$ for every $n\ge N_0$ and every $x$. Define
\begin{equation}
 u(x)=\inf_{n\ge0}S_n\phi(x),\qquad S_0\phi=0.
 \label{eq:positive-inf}
\end{equation}
For $n\ge N_0$ one has $S_n\phi(x)\ge cn>0=S_0\phi(x)$. Hence
\[
 u(x)=\min_{0\le n<N_0}S_n\phi(x),
\]
a minimum of finitely many continuous functions; therefore $u$ is
continuous.

Moreover,
\begin{align*}
\phi(x)+u(Tx)
&=\phi(x)+\inf_{n\ge0}S_n\phi(Tx)\\
&=\inf_{n\ge0}\bigl(\phi(x)+S_n\phi(Tx)\bigr)\\
&=\inf_{m\ge1}S_m\phi(x)\\
&\ge \inf_{m\ge0}S_m\phi(x)=u(x).
\end{align*}
Thus $\phi(x)\ge u(x)-u(Tx)$ for every $x$.
\end{proof}

\begin{remark}
For H\"older regularity, the proposition should be combined with the
regularity theorem available for the quotient class. In particular, for
zooming systems one should invoke the H\"older part of the
Mbarki--Santana theorem rather than infer H\"older regularity merely from
the infimum construction.
\end{remark}

\section{Liv\v{s}ic Transfer}

Uniform fiber contraction already implies periodic lifting; it is not an
additional hypothesis in the present compact setting.

\begin{lemma}[Automatic periodic lifting]
\label{lem:periodic-lifting}
If $p\in M$ satisfies $f^np=p$, then there exists a unique
$\hat p\in N_p$ such that $F^n\hat p=\hat p$.
\end{lemma}

\begin{proof}
Because $\pi$ is continuous and $N$ is compact, the fiber
$N_p=\pi^{-1}(p)$ is compact and hence complete. Since
$f^np=p$, semiconjugacy gives $F^n(N_p)\subset N_p$. For
$\hat x,\hat y\in N_p$,
\[
 d_N(F^n\hat x,F^n\hat y)\le\lambda^n d_N(\hat x,\hat y).
\]
Thus $F^n|_{N_p}$ is a strict contraction of the complete metric space
$N_p$ into itself. The Banach fixed-point theorem gives a unique fixed
point $\hat p\in N_p$.
\end{proof}

\begin{theorem}[Liv\v{s}ic transfer]
\label{thm:livsic}
Assume the reduction theorem and suppose that the quotient admits a
Liv\v{s}ic theorem for the regularity class of the induced potential
$\psi$. Let $\varphi$ satisfy
\[
\sum_{j=0}^{n-1}\varphi(F^j\hat p)=0
\]
whenever $F^n\hat p=\hat p$. Then there exists a quotient transfer $v$
with $\psi=v-v\circ f$, and
\[
 \varphi=U-U\circ F,\qquad U=u+v\circ\pi.
\]
In particular, if $u$ and $v\circ\pi$ are H\"older, then $U$ is
H\"older.
\end{theorem}

\begin{proof}
Write $\varphi=\psi\circ\pi+u-u\circ F$. If $f^np=p$, let
$\hat p\in N_p$ be the fixed point of $F^n$ given by
Lemma~\ref{lem:periodic-lifting}. Then
\begin{align*}
0
&=\sum_{j=0}^{n-1}\varphi(F^j\hat p)\\
&=\sum_{j=0}^{n-1}\psi(f^jp)+u(\hat p)-u(F^n\hat p)\\
&=\sum_{j=0}^{n-1}\psi(f^jp).
\end{align*}
Hence the periodic obstruction vanishes for $\psi$. The quotient
Liv\v{s}ic theorem yields $\psi=v-v\circ f$, and semiconjugacy gives
\[
\varphi=(v-v\circ f)\circ\pi+u-u\circ F
=(u+v\circ\pi)-(u+v\circ\pi)\circ F.
\]
\end{proof}

\begin{corollary}[Liv\v{s}ic from two-sided subcohomology]
\label{cor:two-sided-sc-livsic}
Assume that periodic points of $f$ are dense and that the quotient has
property $\mathrm{SC}$ for both $\psi$ and $-\psi$. Suppose
\[
\int_N\varphi\,d\widetilde\mu=0
\qquad\text{for every }\widetilde\mu\in\M_F(N).
\]
Then there exists $U\in C(N)$ such that $\varphi=U-U\circ F$.
If the quotient subactions and the reduction transfer have H\"older
regularity, then $U$ is H\"older.
\end{corollary}

\begin{proof}
By Corollary~\ref{cor:average-correspondence}, every
$\mu\in\M_f(M)$ satisfies $\int\psi\,d\mu=0$. Property $\mathrm{SC}$ applied to $\psi$ gives a continuous $v_1$ with
\[
v_1-v_1\circ f\le\psi.
\]
Applying $\mathrm{SC}$ to $-\psi$ gives a continuous function $\widetilde v_2$
with
\[
\widetilde v_2-\widetilde v_2\circ f\le-\psi.
\]
Set $v_2=-\widetilde v_2$. Then
\[
\psi\le v_2-v_2\circ f,
\]
and therefore
\[
v_1-v_1\circ f\le\psi\le v_2-v_2\circ f.
\]
Set $w=v_2-v_1$. Then $q:=w-w\circ f\ge0$. If $p$ is periodic of
period $n$, telescoping gives
\[
0=w(p)-w(f^np)=\sum_{j=0}^{n-1}q(f^jp).
\]
Every summand is nonnegative, hence $q=0$ on every periodic orbit. Density
of periodic points and continuity imply $q\equiv0$. Thus the lower and
upper coboundaries coincide, so
$\psi=v_1-v_1\circ f$. The main reduction then yields the required
coboundary on $N$.
\end{proof}

\section{Invariant Measures, Minimizing Measures and Ergodic Optimization}

\begin{theorem}[Bijectivity of invariant lifts]
\label{thm:lift-bijection}
Under the geometric setting above, the projection
\[
\boxed{\pi_*:\M_F(N)\longrightarrow\M_f(M)}
\]
is a bijection. The continuous section supplies existence of invariant lifts,
while uniform fiber contraction makes each invariant lift unique. In
particular every $f$-invariant probability measure has a unique
$F$-invariant lift; no ergodicity assumption is required.
\end{theorem}

\begin{proof}
Let $\mu\in\M_f(M)$. The continuous section gives a probability
$\nu_0=s_*\mu$ satisfying $\pi_*\nu_0=\mu$. Put
\[
\nu_n=\frac1n\sum_{j=0}^{n-1}F_*^j\nu_0.
\]
Every weak-$*$ accumulation point $\widetilde\mu$ of $(\nu_n)$ is
$F$-invariant by the Krylov--Bogolyubov argument. Moreover,
\[
\pi_*F_*^j\nu_0=f_*^j\pi_*\nu_0=f_*^j\mu=\mu,
\]
so $\pi_*\widetilde\mu=\mu$. This proves surjectivity.

For injectivity, suppose
$\pi_*\widetilde\mu_1=\pi_*\widetilde\mu_2=\mu$. Disintegrate over the
common factor:
\[
\widetilde\mu_i=\int_M\widetilde\mu_{i,x}\,d\mu(x),
\qquad \widetilde\mu_{i,x}(N_x)=1
\quad\text{for $\mu$-a.e. }x.
\]
Use the relatively independent joining
\begin{equation}
\kappa=\int_M
(\widetilde\mu_{1,x}\otimes\widetilde\mu_{2,x})\,d\mu(x).
\label{eq:relative-product}
\end{equation}
This is a coupling of $\widetilde\mu_1$ and $\widetilde\mu_2$ supported on
pairs in the same fiber. Since the marginals are invariant,
$(F^n\times F^n)_*\kappa$ is again a coupling of the same two measures.
Therefore
\begin{align*}
W_1(\widetilde\mu_1,\widetilde\mu_2)
&\le \int d_N(F^n\hat x,F^n\hat y)\,d\kappa(\hat x,\hat y)\\
&\le \lambda^n\int d_N(\hat x,\hat y)\,d\kappa\\
&\le \lambda^n\diam(N).
\end{align*}
Letting $n\to\infty$ gives $W_1(\widetilde\mu_1,\widetilde\mu_2)=0$,
so the two measures coincide.
\end{proof}

\begin{remark}
The contraction argument proving uniqueness of an invariant lift uses only the
fiber geometry. The continuous global section is used here as a convenient
device to construct a lift and is essential elsewhere for the explicit
continuous cohomological reduction. Thus the topological section hypothesis
should be regarded primarily as a hypothesis of the reduction mechanism,
rather than of the contraction-based uniqueness argument itself.
\end{remark}

\begin{corollary}[Invariant averages]
\label{cor:average-correspondence}
If $\varphi=\psi\circ\pi+u-u\circ F$ and
$\widetilde\mu\leftrightarrow\mu$ under the bijection above, then
\begin{equation}
\int_N\varphi\,d\widetilde\mu=\int_M\psi\,d\mu.
\label{eq:average-correspondence-v4}
\end{equation}
Consequently $m(\varphi,F)=m(\psi,f)$.
\end{corollary}

\begin{proof}
Invariance gives
$\int(u-u\circ F)\,d\widetilde\mu=0$, while
$\int\psi\circ\pi\,d\widetilde\mu=\int\psi\,d\mu$. The equality of
minimum values follows because $\pi_*$ is bijective on invariant measures.
\end{proof}

Let
\[
\Lambda(\varphi,F)=\left\{\widetilde\mu\in\M_F(N):
\int\varphi\,d\widetilde\mu=m(\varphi,F)\right\},
\]
and define $\Lambda(\psi,f)$ analogously.

\begin{theorem}[Reduction of ergodic optimization]
\label{thm:ergodic-optimization}
The projection restricts to a bijection
\[
\boxed{\pi_*:\Lambda(\varphi,F)\longrightarrow\Lambda(\psi,f).}
\]
Thus uniqueness or non-uniqueness of minimizing measures, and ergodicity of
the corresponding minimizing measure, are preserved by the reduction.
\end{theorem}

\begin{proof}
For corresponding invariant measures the objective functions agree by
\eqref{eq:average-correspondence-v4}, and their minimum values agree by the
preceding corollary. Hence one measure is minimizing if and only if its
projection/lift is minimizing. Bijectivity follows from
Theorem~\ref{thm:lift-bijection}. It remains only to justify the statement
about ergodicity. The map $\pi_*$ is affine. Its inverse is affine as well:
if $\mu=t\mu_1+(1-t)\mu_2$ and $\widetilde\mu_i$ is the unique invariant
lift of $\mu_i$, then $t\widetilde\mu_1+(1-t)\widetilde\mu_2$ is an invariant
lift of $\mu$, hence equals the unique lift of $\mu$. Thus $\pi_*$ is an
affine bijection between the compact convex sets $\M_F(N)$ and $\M_f(M)$.
Affine bijections preserve extreme points, and the extreme points of the
invariant-measure simplex are precisely the ergodic measures. Hence
ergodicity is preserved in both directions.
\end{proof}

\section{Relative Entropy, Pressure and Equilibrium States}\label{sec:pressure}

We first fix the pressure convention used below. For a continuous map
$T:X\to X$ of a compact metric space and a continuous potential $\phi$, set
\begin{equation}
P_T(\phi)
:=\sup_{\mu\in\M_T}
\left(h_\mu(T)+\int\phi\,d\mu\right).
\label{eq:variational-pressure-definition}
\end{equation}
By the variational principle for continuous maps of compact metric spaces,
this quantity agrees with the usual topological pressure defined via
separated (equivalently, spanning) sets; see, for instance, Bowen, Ruelle and Walters
\cite{Bowen,Ruelle,Walters}. Thus no distinction between variational and topological
pressure will be needed below.
When equilibrium states are discussed, we work in the standard finite-pressure
regime; this includes all concrete applications considered below.

For a compact set $K\subset N$, let $s_n(\eps,K)$ denote the maximal
cardinality of an $(n,\eps)$-separated subset of $K$ for the Bowen metric
$d_n^F$. We use the fiberwise relative topological entropy
\begin{equation}
h_{\mathrm{top}}(F\mid\pi)
:=\sup_{x\in M} h_{\mathrm{top}}(F,N_x),
\qquad N_x:=\pi^{-1}(x),
\label{eq:relative-topological-entropy-definition}
\end{equation}
where
\[
h_{\mathrm{top}}(F,K)
:=\lim_{\eps\downarrow0}\limsup_{n\to\infty}
\frac1n\log s_n(\eps,K).
\]
This is the fiberwise entropy quantity entering the version of the
Ledrappier--Walters relative variational inequality used below. We use this
explicit fiberwise formulation throughout; no stronger relative-entropy
notion is required for the pressure argument.

\begin{lemma}[Zero relative topological entropy]
\label{lem:relative-zero}
Under uniform fiber contraction,
\[
\boxed{h_{\mathrm{top}}(F\mid\pi)=0.}
\]
More precisely, for every $\eps>0$, every $x\in M$, and every $n\ge1$,
the maximal cardinality $s_n(\eps,N_x)$ of an $(n,\eps)$-separated subset
of $N_x$ satisfies
\begin{equation}
 s_n(\eps,N_x)\le s_1(\eps,N),
 \label{eq:relative-separated-bound}
\end{equation}
where the right-hand side is finite and independent of $x$ and $n$.
\end{lemma}

\begin{proof}
For $\hat x,\hat y\in N_x$ and $j\ge0$,
\[
d_N(F^j\hat x,F^j\hat y)\le\lambda^j d_N(\hat x,\hat y)
\le d_N(\hat x,\hat y).
\]
Hence the Bowen metric on a fiber is exactly
\[
d_n^F(\hat x,\hat y)
:=\max_{0\le j<n}d_N(F^j\hat x,F^j\hat y)
=d_N(\hat x,\hat y),
\]
because the term $j=0$ occurs in the maximum. Thus every
$(n,\eps)$-separated subset of $N_x$ is already $\eps$-separated for the
original metric. Since $N$ is compact, $s_1(\eps,N)<\infty$, proving
\eqref{eq:relative-separated-bound}. Therefore, uniformly in $x$,
\[
\limsup_{n\to\infty}\frac1n\log s_n(\eps,N_x)=0.
\]
Letting $\eps\downarrow0$ and then taking the supremum over $x$ proves
\eqref{eq:relative-topological-entropy-definition} is zero.
\end{proof}

\begin{lemma}[Relative entropy inequality]
\label{lem:LW-relative}
In the compact semiconjugacy
\[
\pi\circ F=f\circ\pi,
\]
fix $\mu\in\M_f(M)$. The Ledrappier--Walters relative variational principle
gives
\begin{equation}
\sup_{\substack{\widetilde\nu\in\M_F(N)\\ \pi_*\widetilde\nu=\mu}}
 h_{\widetilde\nu}(F)
=
h_\mu(f)+\int_M h_{\mathrm{top}}(F,N_x)\,d\mu(x).
\label{eq:LW-exact-relative}
\end{equation}
Consequently, for every invariant lift $\widetilde\mu$ of $\mu$,
\[
h_\mu(f)\le h_{\widetilde\mu}(F)
\le h_\mu(f)+h_{\mathrm{top}}(F\mid\pi).
\]
Under Lemma~\ref{lem:relative-zero}, every invariant lift therefore satisfies
\[
h_{\widetilde\mu}(F)=h_\mu(f).
\]
\end{lemma}

\begin{proof}
Equation \eqref{eq:LW-exact-relative} is the relative variational principle of
Ledrappier--Walters~\cite{LedrappierWalters}. The lower bound
$h_\mu(f)\le h_{\widetilde\mu}(F)$ is the monotonicity of metric entropy under
factors. For the upper bound, every particular lift is bounded by the
supremum in \eqref{eq:LW-exact-relative}, and
\[
\int_M h_{\mathrm{top}}(F,N_x)\,d\mu(x)
\le h_{\mathrm{top}}(F\mid\pi).
\]
Lemma~\ref{lem:relative-zero} then gives equality of the two metric entropies.
\end{proof}

\begin{theorem}[Entropy and pressure reduction]
\label{thm:pressure}
Let $\widetilde\mu\leftrightarrow\mu$ be corresponding invariant measures.
Then
\begin{equation}
\boxed{h_{\widetilde\mu}(F)=h_\mu(f).}
\label{eq:entropy-equality}
\end{equation}
If $\varphi=\psi\circ\pi+u-u\circ F$, then
\begin{equation}
\boxed{P_F(\varphi)=P_f(\psi).}
\label{eq:pressure-equality}
\end{equation}
Moreover $\widetilde\mu$ is an equilibrium state for $(F,\varphi)$ if and
only if $\mu$ is an equilibrium state for $(f,\psi)$.
\end{theorem}

\begin{proof}
Lemma~\ref{lem:LW-relative}, together with
Lemma~\ref{lem:relative-zero}, gives
\eqref{eq:entropy-equality}. By Corollary~\ref{cor:average-correspondence},
\[
\int_N\varphi\,d\widetilde\mu=\int_M\psi\,d\mu.
\]
Hence the free-energy functionals agree for every corresponding pair:
\begin{equation}
h_{\widetilde\mu}(F)+\int_N\varphi\,d\widetilde\mu
=
h_\mu(f)+\int_M\psi\,d\mu.
\label{eq:free-energy-correspondence}
\end{equation}
Since Theorem~\ref{thm:lift-bijection} gives a bijection
$\pi_*:\M_F(N)\to\M_f(M)$, taking suprema in
\eqref{eq:free-energy-correspondence} yields
\[
P_F(\varphi)=P_f(\psi).
\]
The same pointwise identity of free-energy functionals, together with the
bijection of invariant measures, proves the equilibrium-state
correspondence.
\end{proof}

\section{Weak Gibbs Transfer under Quantitative Bowen-Mass Comparison}\label{sec:weak-gibbs}

For $T:X\to X$ write
\[
B_T(x,\eps,n)=\{y:d(T^jy,T^jx)<\eps,\ 0\le j<n\}.
\]
Throughout this section we assume that the pressures appearing in the weak
Gibbs expressions are finite. Let $G_T\subset X$ be a prescribed Gibbs set.
We use the uniform-on-$G_T$ convention for weak Gibbs estimates: the
subexponential constants below depend on the radius and on $n$, but not on
the center $x\in G_T$. We say that a probability measure $\nu$ is
\emph{uniform weak Gibbs for $(T,\phi)$ on $G_T$} if for every
sufficiently small $\eps>0$ there are numbers $K_n(\eps)\ge1$ such that
\begin{equation}
\lim_{n\to\infty}\frac1n\log K_n(\eps)=0
\label{eq:weak-gibbs-subexp}
\end{equation}
and, for every $x\in G_T$ and every $n\ge1$,
\begin{equation}
K_n(\eps)^{-1}
\le
\frac{\nu(B_T(x,\eps,n))}
{\exp(S_n\phi(x)-nP_T(\phi))}
\le K_n(\eps).
\label{eq:weak-gibbs-def}
\end{equation}
We abbreviate this uniform-on-$G_T$ property simply as ``weak Gibbs'' in the
statements below. This convention is stronger than pointwise variants in
which the subexponential error is allowed to depend on the center.

Zero relative entropy alone does not compare the measures of individual
Bowen balls. The appropriate additional input is a subexponential
comparison of Bowen-ball masses for a \emph{specific corresponding pair}
of invariant measures. The condition below should be read as a sufficient
transfer hypothesis. We do not claim that it is necessary, nor that it follows
from uniform fiber contraction. The latter controls the diameter of the
iterated image of each individual fiber, but it gives no quantitative lower
bound on the conditional mass assigned to a small fiber neighborhood. This
is precisely the additional information needed to compare the masses of
Bowen balls in the extension and in the quotient.

\begin{definition}[Forward subexponential Bowen-mass comparability]
\label{def:forward-bowen-comparison}
Let $\widetilde\mu\in\M_F(N)$ and $\mu=\pi_*\widetilde\mu$. Let
$G_F\subset N$ and $G_f\subset M$ satisfy $\pi(G_F)\subset G_f$. We say
that $(\widetilde\mu,\mu)$ satisfies \emph{forward subexponential
Bowen-mass comparability on $(G_F,G_f)$} if, for every sufficiently small
$\eps>0$, there exist radii
$\eps_-(\eps),\eps_+(\eps)>0$ tending to zero with $\eps$ and numbers
$a_n^{+}(\eps)\ge0$ such that
\[
\lim_{n\to\infty}\frac{a_n^{+}(\eps)}n=0
\]
and, for every $\hat x\in G_F$, $x=\pi(\hat x)$, and every $n\ge1$,
\begin{equation}
e^{-a_n^{+}(\eps)}\mu(B_f(x,\eps_-,n))
\le \widetilde\mu(B_F(\hat x,\eps,n))
\le e^{a_n^{+}(\eps)}\mu(B_f(x,\eps_+,n)).
\label{eq:bowen-comp-forward}
\end{equation}
\end{definition}

\begin{definition}[Backward subexponential Bowen-mass comparability]
\label{def:backward-bowen-comparison}
Let $\widetilde\mu\in\M_F(N)$ and $\mu=\pi_*\widetilde\mu$. Let
$G_F\subset N$ and $G_f\subset M$ be such that every $x\in G_f$ admits
at least one lift $\hat x\in G_F$ with $\pi(\hat x)=x$. We say that
$(\widetilde\mu,\mu)$ satisfies \emph{backward subexponential Bowen-mass
comparability on $(G_F,G_f)$} if, for every sufficiently small $\eps>0$,
there exist radii $\eps_-(\eps),\eps_+(\eps)>0$ tending to zero with
$\eps$ and numbers $a_n^{-}(\eps)\ge0$ such that
\[
\lim_{n\to\infty}\frac{a_n^{-}(\eps)}n=0
\]
and, for every $x\in G_f$, for a prescribed lift $\hat x\in G_F$ with
$\pi(\hat x)=x$, and every $n\ge1$,
\begin{equation}
e^{-a_n^{-}(\eps)}\widetilde\mu(B_F(\hat x,\eps_-,n))
\le \mu(B_f(x,\eps,n))
\le e^{a_n^{-}(\eps)}\widetilde\mu(B_F(\hat x,\eps_+,n)).
\label{eq:bowen-comp-backward}
\end{equation}
\end{definition}

The logarithmic formulation above makes explicit that only an $o(n)$ error
in Bowen-ball masses is required. Equivalently one may write
$A_n^{\pm}(\eps)=e^{a_n^{\pm}(\eps)}$ with
$n^{-1}\log A_n^{\pm}(\eps)\to0$.

\begin{lemma}[Stability under radius reparametrization]
\label{lem:radius-reparametrization}
Let $\nu$ be weak Gibbs for $(T,\phi)$ in the sense of
\eqref{eq:weak-gibbs-def}. Suppose that, for every sufficiently small
$\eps>0$, two radii $\delta_-(\eps),\delta_+(\eps)>0$ satisfy
\[
\delta_\pm(\eps)\longrightarrow0\qquad(\eps\to0).
\]
Then the weak Gibbs estimates may be evaluated at
$\delta_-(\eps)$ and $\delta_+(\eps)$ without changing the subexponential
character of the Gibbs constants. More precisely, if
$K_n(\delta_\pm(\eps))$ are the corresponding Gibbs constants, then
\[
\frac1n\log
\max\{K_n(\delta_-(\eps)),K_n(\delta_+(\eps))\}\longrightarrow0.
\]
The same conclusion remains valid after finitely many further radius
reparametrizations, provided all resulting radii tend to zero with the
target radius.
\end{lemma}

\begin{proof}
For each fixed sufficiently small target radius $\eps$, the radii
$\delta_\pm(\eps)$ are themselves fixed positive radii. Hence the weak Gibbs
definition applies to each of them, and both corresponding logarithmic
constants are $o(n)$. The maximum of finitely many $o(n)$ sequences is again
$o(n)$.
\end{proof}

\begin{theorem}[Forward weak Gibbs transfer]
\label{thm:weak-gibbs-forward}
Assume
\[
\varphi=\psi\circ\pi+u-u\circ F,
\]
and let $\widetilde\mu\leftrightarrow\mu$ be corresponding invariant
measures. Let $G_F\subset N$ and $G_f\subset M$ satisfy
$\pi(G_F)\subset G_f$. If $(\widetilde\mu,\mu)$ satisfies forward
subexponential Bowen-mass comparability on $(G_F,G_f)$ and $\mu$ is weak
Gibbs for $(f,\psi)$ on $G_f$, then $\widetilde\mu$ is weak Gibbs for
$(F,\varphi)$ on $G_F$.
\end{theorem}

\begin{proof}
For every $\hat x$ and every $n$,
\begin{equation}
S_n\varphi(\hat x)
=S_n\psi(\pi\hat x)+u(\hat x)-u(F^n\hat x),
\label{eq:birkhoff-coboundary}
\end{equation}
so
\begin{equation}
\bigl|S_n\varphi(\hat x)-S_n\psi(\pi\hat x)\bigr|
\le2\|u\|_\infty.
\label{eq:birkhoff-bounded-error}
\end{equation}
Also $P_F(\varphi)=P_f(\psi)$ by
Theorem~\ref{thm:pressure}. Fix $\hat x\in G_F$ and put
$x=\pi(\hat x)\in G_f$. By \eqref{eq:bowen-comp-forward} and the weak Gibbs
bounds for $\mu$ at the radii $\eps_-$ and $\eps_+$,
\begin{align*}
\widetilde\mu(B_F(\hat x,\eps,n))
&\le e^{a_n^+(\eps)}K_n(\eps_+)
 \exp(S_n\psi(x)-nP_f(\psi))\\
&\le e^{a_n^+(\eps)}K_n(\eps_+)e^{2\|u\|_\infty}
 \exp(S_n\varphi(\hat x)-nP_F(\varphi)),
\end{align*}
and analogously
\begin{align*}
\widetilde\mu(B_F(\hat x,\eps,n))
&\ge e^{-a_n^+(\eps)}K_n(\eps_-)^{-1}
 \exp(S_n\psi(x)-nP_f(\psi))\\
&\ge e^{-a_n^+(\eps)}K_n(\eps_-)^{-1}e^{-2\|u\|_\infty}
 \exp(S_n\varphi(\hat x)-nP_F(\varphi)).
\end{align*}
Thus one may take, for instance,
\[
\widetilde K_n(\eps)
=e^{a_n^+(\eps)+2\|u\|_\infty}
\max\{K_n(\eps_-),K_n(\eps_+)\}.
\]
Then $n^{-1}\log\widetilde K_n(\eps)\to0$, proving the claim.
\end{proof}

\begin{theorem}[Backward weak Gibbs transfer]
\label{thm:weak-gibbs-backward}
Assume
$\varphi=\psi\circ\pi+u-u\circ F$, and let
$\widetilde\mu\leftrightarrow\mu$ be corresponding invariant measures.
Let $G_F\subset N$ and $G_f\subset M$ be such that every $x\in G_f$ has
a prescribed lift $\hat x\in G_F$. If $(\widetilde\mu,\mu)$ satisfies
backward subexponential Bowen-mass comparability on $(G_F,G_f)$ and
$\widetilde\mu$ is weak Gibbs for $(F,\varphi)$ on $G_F$, then $\mu$ is
weak Gibbs for $(f,\psi)$ on $G_f$.
\end{theorem}

\begin{proof}
Fix $x\in G_f$ and its prescribed lift $\hat x\in G_F$. Apply the weak
Gibbs bounds for $\widetilde\mu$ at the radii $\eps_-$ and $\eps_+$ and
then use \eqref{eq:bowen-comp-backward}. Combining these inequalities with
\eqref{eq:birkhoff-bounded-error} and
$P_F(\varphi)=P_f(\psi)$ yields
\begin{align*}
\mu(B_f(x,\eps,n))
&\le e^{a_n^-(\eps)}K_n(\eps_+)e^{2\|u\|_\infty}
\exp(S_n\psi(x)-nP_f(\psi)),\\
\mu(B_f(x,\eps,n))
&\ge e^{-a_n^-(\eps)}K_n(\eps_-)^{-1}e^{-2\|u\|_\infty}
\exp(S_n\psi(x)-nP_f(\psi)).
\end{align*}
The corresponding Gibbs constant is again subexponential, since
$a_n^-(\eps)=o(n)$ and \eqref{eq:weak-gibbs-subexp} holds for
$\widetilde\mu$.
\end{proof}

\begin{corollary}[Two-sided weak Gibbs equivalence]
\label{thm:weak-gibbs}
Assume $\varphi=\psi\circ\pi+u-u\circ F$. Let
$\widetilde\mu\leftrightarrow\mu$, and let $G_F,G_f$ be Gibbs sets for
which both forward and backward subexponential Bowen-mass comparability
hold. Then
\[
\mu\text{ is weak Gibbs for }(f,\psi)\text{ on }G_f
\quad\Longleftrightarrow\quad
\widetilde\mu\text{ is weak Gibbs for }(F,\varphi)\text{ on }G_F.
\]
\end{corollary}

\begin{proof}
The forward implication is Theorem~\ref{thm:weak-gibbs-forward}, and the
reverse implication is Theorem~\ref{thm:weak-gibbs-backward}. The common
pressure identity and the bounded coboundary error are already contained in
those proofs. Hence no additional assumption is required beyond the two
Bowen-mass comparisons.
\end{proof}

\begin{remark}[What is automatic and what is not]
\label{rem:bowen-comparison-not-automatic}
The subexponential Bowen-mass comparison above is strictly stronger than
zero relative topological entropy and is not implied by the pressure
identity. Uniform continuity of $\pi$ gives one geometric inclusion for
Bowen balls: for every $\eta>0$ there exists $\eps>0$ such that
\[
\pi(B_F(\hat x,\eps,n))\subset B_f(\pi\hat x,\eta,n)
\quad\text{for every }\hat x\in N\text{ and }n\ge1.
\]
Consequently,
\[
\widetilde\mu(B_F(\hat x,\eps,n))
\le \mu(B_f(\pi\hat x,\eta,n)).
\]
The reverse comparison is not a consequence of fiber contraction alone:
contraction controls pairs of points that already lie in the same fiber,
whereas a base Bowen ball contains points in distinct fibers. Therefore
forward/backward Bowen-mass comparability is an additional quantitative
geometric/measure-theoretic hypothesis that must be verified in each
application.
\end{remark}


\subsection{Concrete geometric criteria for Bowen-mass comparability}

The abstract comparison hypotheses above can be verified from more primitive
geometric information in several important classes. We record three useful
criteria. They are stated only as sufficient conditions; none of them is
claimed to follow from uniform fiber contraction alone.

Whenever radius moduli are introduced below, we choose them monotonically
(after shrinking them if necessary). In particular, the lower radius moduli
are chosen cofinal at the origin: for every sufficiently small target radius
$\eta>0$ there is a sufficiently small $\rho>0$ for which
$\eta\le\delta_-(\rho)$. This harmless convention makes the backward
reparametrizations used below explicit.

\begin{proposition}[Invariant-graph criterion]
\label{prop:invariant-graph-bowen}
Assume that there exists a continuous invariant section
\[
s:M\longrightarrow N,
\qquad \pi\circ s=\mathrm{id}_M,
\qquad F\circ s=s\circ f.
\]
Let $\mu\in\M_f(M)$ and let $\widetilde\mu$ be its unique invariant lift.
Then $\widetilde\mu=s_*\mu$. If $G_f\subset M$ and $G_F=s(G_f)$, then
$(\widetilde\mu,\mu)$ satisfies both forward and backward subexponential
Bowen-mass comparability on $(G_F,G_f)$ with
\[
a_n^+(\eps)=a_n^-(\eps)=0.
\]
More precisely, for every sufficiently small $\eps>0$ there are radii
$\delta_-(\eps),\delta_+(\eps)>0$, tending to zero with $\eps$, such that for
every $x\in G_f$, every $n\ge1$, and $\hat x=s(x)$,
\[
\mu(B_f(x,\delta_-,n))
\le
\widetilde\mu(B_F(\hat x,\eps,n))
\le
\mu(B_f(x,\delta_+,n)).
\]
\end{proposition}

\begin{proof}
Since $\mu$ is $f$-invariant and $F\circ s=s\circ f$,
\[
F_*s_*\mu=s_*f_*\mu=s_*\mu.
\]
Moreover $\pi_*s_*\mu=\mu$. By Theorem~\ref{thm:lift-bijection}, the invariant
lift is unique, hence $\widetilde\mu=s_*\mu$.

Fix $\eps>0$. Since $M$ is compact and $s$ is continuous, $s$ is uniformly
continuous. Choose $\delta_->0$ such that
\[
d_M(x,y)<\delta_-
\quad\Longrightarrow\quad
d_N(s(x),s(y))<\eps.
\]
If $y\in B_f(x,\delta_-,n)$, then for every $0\le j<n$,
\[
d_N(F^js(x),F^js(y))
=d_N(s(f^jx),s(f^jy))<\eps,
\]
so
\[
s(B_f(x,\delta_-,n))\subset B_F(s(x),\eps,n)\cap s(M).
\]
Conversely, $\pi|_{s(M)}:s(M)\to M$ is continuous and $s(M)$ is compact;
therefore it is uniformly continuous. Choose $\delta_+>0$ so that
\[
d_N(s(x),s(y))<\eps
\quad\Longrightarrow\quad
d_M(x,y)<\delta_+.
\]
Then
\[
B_F(s(x),\eps,n)\cap s(M)
\subset s(B_f(x,\delta_+,n)).
\]
Since $\widetilde\mu=s_*\mu$ is supported on $s(M)$, taking measures gives the
claimed inequalities. Forward comparability follows immediately. Backward
comparability follows by applying the same inclusions after reparametrizing
the radii with the uniform-continuity moduli of $s$ and
$\pi|_{s(M)}=s^{-1}$. No multiplicative mass error is introduced, so the
logarithmic errors are identically zero.
\end{proof}

The next criterion isolates the measure-theoretic nondegeneracy needed to
complement the geometric inclusion provided by the skew-product structure.

\begin{proposition}[Uniform skew-product criterion]
\label{prop:uniform-skew-bowen}
Assume that $N=M\times Y$ is endowed with a metric uniformly equivalent to a
product metric and
\[
F(x,y)=\bigl(f(x),G_x(y)\bigr).
\]
Assume uniform fiber contraction: there is $0<\lambda<1$ such that
\[
d_Y(G_x(y),G_x(y'))\le\lambda d_Y(y,y')
\]
for all admissible points. Assume also uniform transverse continuity: there is
a modulus $\omega:[0,\infty)\to[0,\infty)$ with $\omega(t)\to0$ as $t\to0$
and
\[
d_Y(G_x(y),G_{x'}(y))\le \omega(d_M(x,x'))
\]
for all admissible points.

Let
\[
\widetilde\mu=\int_M\widetilde\mu_x\,d\mu(x)
\]
be the disintegration over $\pi(x,y)=x$. Let $G_F\subset N$ and
$G_f=\pi(G_F)$. Assume the following \emph{averaged subexponential fiber
nondegeneracy}: for every sufficiently small fiber radius $r>0$ and every
sufficiently small base Bowen radius $\delta>0$ there are numbers
$q_n(r,\delta)\ge0$ such that
\[
\frac{q_n(r,\delta)}{n}\longrightarrow0
\qquad (n\to\infty),
\]
and, uniformly for $(x,y)\in G_F$ and every $n\ge1$,
\begin{equation}
\label{eq:averaged-lower-fiber-ball-masses}
\int_{B_f(x,\delta,n)}
\widetilde\mu_{x'}(B_Y(y,r))\,d\mu(x')
\ge e^{-q_n(r,\delta)}\,\mu(B_f(x,\delta,n)).
\end{equation}
Then $(\widetilde\mu,\mu)$ satisfies forward and backward subexponential
Bowen-mass comparability on $(G_F,G_f)$.
\end{proposition}

\begin{proof}
Fix a sufficiently small Bowen radius $\eps>0$. Uniform equivalence of the
ambient metric with a product metric provides product radii tending to zero
with $\eps$; thus it suffices to establish the desired inclusions in product
coordinates with radii depending only on $\eps$.

Let $(x,y),(x',y')\in M\times Y$, and write their fiber coordinates after $j$
iterates as $y_j,y_j'$. Iterating the contraction/transverse-continuity bounds
gives
\[
d_Y(y_j,y_j')
\le \lambda^j d_Y(y,y')
+\sum_{k=0}^{j-1}\lambda^{j-1-k}
\omega\bigl(d_M(f^kx,f^kx')\bigr).
\]
Hence, if $x'\in B_f(x,\delta,n)$ and $d_Y(y,y')<r$, then
\[
d_Y(y_j,y_j')
\le r+\frac{\omega(\delta)}{1-\lambda},
\qquad 0\le j<n.
\]
Choose first $r=r(\eps)>0$ and then $\delta_-(\eps)>0$ so that the right-hand
side lies inside the fiber component of the chosen Bowen radius. Together with
the base component of the product metric this gives
\begin{equation}
\label{eq:skew-product-lower-inclusion}
\bigl\{(x',y'):\ x'\in B_f(x,\delta_-,n),\
 y'\in B_Y(y,r)\bigr\}
\subset B_F((x,y),\eps,n).
\end{equation}
By uniform continuity of $\pi$, for some $\delta_+(\eps)\to0$,
\begin{equation}
\label{eq:skew-product-upper-inclusion}
B_F((x,y),\eps,n)
\subset \pi^{-1}(B_f(x,\delta_+,n)).
\end{equation}
Disintegrating and using
\eqref{eq:averaged-lower-fiber-ball-masses} with
$r=r(\eps)$ and $\delta=\delta_-(\eps)$ in the lower inclusion gives
\[
\widetilde\mu(B_F((x,y),\eps,n))
\ge e^{-q_n(r(\eps),\delta_-(\eps))}
\mu(B_f(x,\delta_-(\eps),n)),
\]
whereas the upper inclusion gives
\[
\widetilde\mu(B_F((x,y),\eps,n))
\le \mu(B_f(x,\delta_+(\eps),n)).
\]
Thus forward comparability holds with
\[
a_n^+(\eps)=q_n(r(\eps),\delta_-(\eps))=o(n).
\]

For the backward comparison, fix a small base radius $\eta>0$. Choose
$\rho_-(\eta)>0$ so small that
$\delta_+(\rho_-(\eta))\le\eta$. Applying the upper inequality above at radius
$\rho_-(\eta)$ yields
\[
\widetilde\mu(B_F((x,y),\rho_-(\eta),n))
\le \mu(B_f(x,\eta,n)).
\]
Next choose $\rho_+(\eta)>0$ so that
$\eta\le\delta_-(\rho_+(\eta))$. Applying the lower inequality at radius
$\rho_+(\eta)$ and using monotonicity of base Bowen balls gives
\[
\mu(B_f(x,\eta,n))
\le e^{q_n(r(\rho_+(\eta)),\delta_-(\rho_+(\eta)))}
\widetilde\mu(B_F((x,y),\rho_+(\eta),n)).
\]
Therefore backward comparability holds with
\[
a_n^-(\eta)
=q_n(r(\rho_+(\eta)),\delta_-(\rho_+(\eta)))=o(n).
\]
Hence the two-sided subexponential Bowen-mass comparison follows.
\end{proof}

\begin{remark}
The pointwise lower bound
$\widetilde\mu_{x'}(B_Y(y,r))\ge c_r>0$ used in more uniform settings is a
special case of \eqref{eq:averaged-lower-fiber-ball-masses}, corresponding to
$q_n(r,\delta)\equiv |\log c_r|$. The present formulation allows spatial
nonuniformity, but the subexponential envelope must be uniform over the Gibbs
set. Mere pointwise convergence of center-dependent errors is not enough for
the uniform weak Gibbs definition used above.
\end{remark}

\begin{proposition}[Local-product chart criterion]
\label{prop:local-product-bowen}
Let $G_F\subset N$. Assume that for every $\hat x\in G_F$ there are
neighborhoods $U_{\hat x}^{\mathrm{base}}\subset M$,
$U_{\hat x}^{\mathrm{fib}}$ in a fixed local fiber model, a neighborhood
$R_{\hat x}\subset N$, and a homeomorphism
\[
\Theta_{\hat x}:
U_{\hat x}^{\mathrm{base}}\times U_{\hat x}^{\mathrm{fib}}
\longrightarrow R_{\hat x}
\]
such that
\[
\pi\circ\Theta_{\hat x}(z,\xi)=z
\]
whenever the expression is defined.

Assume that for every sufficiently small $\eps>0$ there are a base radius
$\delta_-(\eps)>0$, a fiber radius $r(\eps)>0$, local product boxes
$Q_n(\hat x,\eps)\subset R_{\hat x}$, and a radius modulus
$\rho(\eps)\to0$ such that, for every $\hat x\in G_F$ and every $n\ge1$,
\begin{equation}
\label{eq:local-product-chart-lower}
\Theta_{\hat x}\!\left(
B_f(\pi\hat x,\delta_-(\eps),n)\times
B^{\mathrm{fib}}(\hat x,r(\eps))
\right)
\subset
Q_n(\hat x,\eps)
\subset
B_F(\hat x,\rho(\eps),n).
\end{equation}
Assume moreover a lower product nondegeneracy estimate: for every sufficiently
small $\eps>0$ there are numbers $b_n(\eps)\ge0$ with $b_n(\eps)=o(n)$ and a
constant $c_\eps>0$ such that
\begin{align}
\label{eq:local-product-chart-mass}
&\widetilde\mu\!\left(
\Theta_{\hat x}\!\left(
B_f(\pi\hat x,\delta_-(\eps),n)\times
B^{\mathrm{fib}}(\hat x,r(\eps))
\right)\right)
\nonumber\\
&\hspace{2cm}\ge
c_\eps e^{-b_n(\eps)}
\mu(B_f(\pi\hat x,\delta_-(\eps),n))
\end{align}
for every $\hat x\in G_F$ and every $n\ge1$.

Then $(\widetilde\mu,\mu)$ satisfies both forward and backward
subexponential Bowen-mass comparability on $(G_F,\pi(G_F))$. More precisely,
for every small target radius one may take logarithmic errors equal to
\[
\widehat b_n(\eps)+O_\eps(1),
\]
where $\widehat b_n(\eps)$ is the maximum of the finitely many
$b_n(\rho_j(\eps))$ generated by the radius reparametrizations used in the
forward and backward comparisons. Hence $\widehat b_n(\eps)=o(n)$. If the
lower product distortion is uniformly bounded in $n$, then
$a_n^\pm(\eps)=O_\eps(1)$.
\end{proposition}

\begin{proof}
The first inclusion in \eqref{eq:local-product-chart-lower} and the lower mass
estimate \eqref{eq:local-product-chart-mass} give
\[
\widetilde\mu(B_F(\hat x,\rho(\eps),n))
\ge
c_\eps e^{-b_n(\eps)}
\mu(B_f(\pi\hat x,\delta_-(\eps),n)).
\]
On the other hand, uniform continuity of $\pi$ gives, for every sufficiently
small Bowen radius $\rho>0$, a base radius $\eta(\rho)\to0$ such that
\[
B_F(\hat x,\rho,n)
\subset
\pi^{-1}(B_f(\pi\hat x,\eta(\rho),n)),
\]
and therefore
\[
\widetilde\mu(B_F(\hat x,\rho,n))
\le
\mu(B_f(\pi\hat x,\eta(\rho),n)).
\]
We first make the forward comparison explicit. Given a small target Bowen
radius $\tau>0$, choose a parameter $\eps_\tau>0$ so that
$\rho(\eps_\tau)\le\tau$. Then monotonicity of Bowen balls gives
\[
\widetilde\mu(B_F(\hat x,\tau,n))
\ge c_{\eps_\tau}e^{-b_n(\eps_\tau)}
\mu(B_f(\pi\hat x,\delta_-(\eps_\tau),n)).
\]
Together with the projection estimate at radius $\tau$, this is forward
comparability with logarithmic error
$b_n(\eps_\tau)+|\log c_{\eps_\tau}|=o(n)$.

For backward comparability, fix a small base radius $\eta>0$. By the
monotone radius convention, choose a Bowen radius $\tau_-(\eta)>0$ such that
$\eta(\tau_-(\eta))\le\eta$. The projection estimate then gives
\[
\widetilde\mu(B_F(\hat x,\tau_-(\eta),n))
\le\mu(B_f(\pi\hat x,\eta,n)).
\]
Next, by cofinality of $\delta_-$ at the origin, choose a parameter
$\eps_+(\eta)>0$ so that
\[
\eta\le\delta_-(\eps_+(\eta)).
\]
The lower product estimate and monotonicity of base Bowen balls yield
\[
\mu(B_f(\pi\hat x,\eta,n))
\le c_{\eps_+(\eta)}^{-1}e^{b_n(\eps_+(\eta))}
\widetilde\mu(B_F(\hat x,\rho(\eps_+(\eta)),n)).
\]
Thus the backward logarithmic error is
$b_n(\eps_+(\eta))+|\log c_{\eps_+(\eta)}|=o(n)$. Only finitely many
radius reparametrizations are used for each target radius, so one may collect
their errors into
\[
\widehat b_n(\eps)+O_\eps(1),
\qquad
\widehat b_n(\eps)=\max_j b_n(\rho_j(\eps))=o(n).
\]
This proves the two-sided comparison.
\end{proof}


\section{Sequential Gibbs Measures and Fiber-Contraction}\label{sec:sequential-gibbs}

The weak Gibbs results above require estimates at every time. In many
non-uniformly hyperbolic situations, however, Gibbs estimates are available
only along a point-dependent subsequence of times. This motivates a
sequential version of the transfer mechanism. The formulation below is
adapted to Bowen balls and is deliberately separated from the symbolic
cylinder definition of Ferreira--Oliveira \cite{FerreiraOliveira}; the
relation with that definition is discussed at the end of the section.

\begin{definition}[Bowen sequential Gibbs]
\label{def:bowen-sequential-gibbs}
Let $T:X\to X$ be continuous, let $\phi\in C(X)$, let $\nu$ be a probability
measure, and let $G\subset X$ be measurable with $\nu(G)=1$. We say that
$\nu$ is \emph{Bowen sequential Gibbs for $(T,\phi)$ on $G$} if there exist
$P\in\mathbb R$ and, for every $x\in G$, an increasing sequence
$n_1(x)<n_2(x)<\cdots$, chosen independently of the Bowen radius, such that
the following uniform property holds: for every sufficiently small
$\eps>0$ there exists $K(\eps)\ge1$, independent of $x$, $i$, and $j$,
such that, for every $x\in G$, every $i\ge1$, and every
$0\le j<n_i(x)$,
\begin{equation}
 K(\eps)^{-1}
 \le
 \frac{\nu\bigl(B_T(T^jx,\eps,n_i(x)-j)\bigr)}
 {\exp\bigl(S_{n_i(x)-j}\phi(T^jx)-(n_i(x)-j)P\bigr)}
 \le K(\eps).
 \label{eq:bowen-sequential-gibbs}
\end{equation}
Any such sequence is called an \emph{admissible Bowen Gibbs-time sequence}
for $x$, and $P$ is called the \emph{sequential Gibbs pressure} associated
with these estimates.
\end{definition}

\begin{remark}[Uniformity convention]
The constant $K(\eps)$ in Definition~\ref{def:bowen-sequential-gibbs} is
uniform on the full-measure Gibbs set $G$. Only the admissible sequence
$n_i(x)$ is allowed to depend on the center. This matches the uniform Gibbs
constant in the classical symbolic sequential Gibbs formulation, while the
Bowen radii are permitted to vary through the radius reparametrizations used
below.
\end{remark}

\begin{remark}[Admissible versus maximal Gibbs times]
The preceding definition records one admissible sequence along which all
suffix estimates hold. We do not claim that this sequence is maximal among
all admissible times. Accordingly, the transfer results below prove that a
prescribed admissible Gibbs-time sequence remains admissible after passage
between quotient and extension. This is weaker than identifying the maximal
set of all Gibbs times.
\end{remark}

The tail requirement in \eqref{eq:bowen-sequential-gibbs} mirrors the
standard symbolic definition of sequential Gibbs measures, in which the
same estimate is required for every suffix of a Gibbs cylinder. We do not
identify $P$ a priori with full topological pressure: in the symbolic theory
it is naturally associated with the pressure carried by the Gibbs set.

\begin{definition}[Blockwise Bowen sequential weak Gibbs]
\label{def:bowen-sequential-weak}
With the notation above, $\nu$ is \emph{blockwise Bowen sequential weak
Gibbs} with sequential pressure $P$ if, for every $x\in G$, there is an
increasing sequence $n_i(x)\to\infty$, chosen independently of the Bowen
radius, such that for every sufficiently small $\eps>0$ there are constants
$K_i(x,\eps)\ge1$ for which \eqref{eq:bowen-sequential-gibbs} holds with
$K(\eps)$ replaced by $K_i(x,\eps)$, uniformly for all
$0\le j<n_i(x)$, and
\begin{equation}
 \frac{1}{n_i(x)}\log K_i(x,\eps)\longrightarrow0.
 \label{eq:sequential-weak-subexp}
\end{equation}
\end{definition}

\begin{remark}[Normalization of the weak error]
The error in Definition~\ref{def:bowen-sequential-weak} is normalized by the
full Gibbs-block length $n_i(x)$, not by the suffix length $n_i(x)-j$. Thus the
definition is intentionally weaker on very short suffixes. The adjective
\emph{blockwise} will be used to avoid confusing this convention with stronger
sequential weak Gibbs notions in which the error is required to be
subexponential relative to every suffix length.
\end{remark}

\begin{definition}[Forward sequential Bowen-mass comparability]
\label{def:sequential-forward-comparison}
Let $\widetilde\mu\leftrightarrow\mu$ be corresponding invariant measures,
let $G_f\subset M$ have full $\mu$-measure, and suppose that an admissible
sequence $n_i(x)$ is fixed for every $x\in G_f$. We say that \emph{forward
sequential Bowen-mass comparability} holds if there is a full
$\widetilde\mu$-measure set $G_F\subset\pi^{-1}(G_f)$ such that, for every
sufficiently small $\eps>0$, there are radii
$\eps_-(\eps),\eps_+(\eps)>0$ tending to zero with $\eps$ and numbers
$a_i(\hat x,\eps)\ge0$ satisfying
\begin{equation}
 \frac{a_i(\hat x,\eps)}{n_i(\pi\hat x)}\longrightarrow0
 \qquad\text{for every }\hat x\in G_F,
 \label{eq:seq-mass-error-subexp}
\end{equation}
such that, writing $x=\pi\hat x$ and $m=n_i(x)-j$, for every
$0\le j<n_i(x)$,
\begin{align}
 e^{-a_i(\hat x,\eps)}
 \mu\bigl(B_f(f^jx,\eps_-,m)\bigr)
 &\le
 \widetilde\mu\bigl(B_F(F^j\hat x,\eps,m)\bigr)
 \nonumber\\
 &\le
 e^{a_i(\hat x,\eps)}
 \mu\bigl(B_f(f^jx,\eps_+,m)\bigr).
 \label{eq:seq-forward-comparison}
\end{align}
If, for each fixed $\eps$, the errors $a_i(\hat x,\eps)$ are bounded
uniformly in $\hat x$ and $i$, we say that the comparison is
\emph{uniformly bounded}.
\end{definition}

\begin{definition}[Backward sequential Bowen-mass comparability]
\label{def:sequential-backward-comparison}
Let $G_f\subset M$ have full $\mu$-measure. Assume that a measurable
prescribed lift
\[
 \ell:G_f\to N,
 \qquad \pi\circ\ell=\operatorname{id}_{G_f},
\]
is fixed, with $\ell(G_f)$ contained in a full $\widetilde\mu$-measure set
$G_F$ on which an admissible Gibbs-time sequence $n_i(\hat x)$ is defined.
Set
\[
 n_i(x):=n_i(\ell(x)).
\]
We say that \emph{backward sequential Bowen-mass comparability along $\ell$}
holds if, for every sufficiently small $\eps>0$, there are radii
$\eps_-(\eps),\eps_+(\eps)>0$ tending to zero with $\eps$ and numbers
$b_i(x,\eps)\ge0$ satisfying
\[
 \frac{b_i(x,\eps)}{n_i(x)}\longrightarrow0
\]
such that, for every $x\in G_f$, every $0\le j<n_i(x)$, and
$m=n_i(x)-j$,
\begin{align}
 e^{-b_i(x,\eps)}
 \widetilde\mu\bigl(B_F(F^j\ell(x),\eps_-,m)\bigr)
 &\le \mu\bigl(B_f(f^jx,\eps,m)\bigr)
 \nonumber\\
 &\le e^{b_i(x,\eps)}
 \widetilde\mu\bigl(B_F(F^j\ell(x),\eps_+,m)\bigr).
 \label{eq:seq-backward-comparison}
\end{align}
If the $b_i(x,\eps)$ are uniformly bounded in $x$ and $i$ for fixed
$\eps$, we call the comparison \emph{uniformly bounded}.
\end{definition}

\begin{remark}[Sequential synchronization along fibers]
A prescribed lift is one convenient way to make the backward Gibbs times a
well-defined function of the base point. Alternatively, one may replace it by
a \emph{fiberwise synchronization hypothesis}: on a full set,
$\pi(\hat x)=\pi(\hat y)$ implies that the chosen admissible Gibbs-time
sequences coincide. Without one of these devices, contraction alone does not
force Gibbs-time sequences chosen upstairs to descend canonically to the
quotient.
\end{remark}

\begin{theorem}[Forward blockwise sequential weak Gibbs transfer]
\label{thm:sequential-weak-forward}
Assume
\[
 \varphi=\psi\circ\pi+u-u\circ F,
\]
and let $\widetilde\mu\leftrightarrow\mu$. Suppose $\mu$ is blockwise Bowen
sequential weak Gibbs for $(f,\psi)$ with sequential pressure $P$ and a fixed
admissible sequence $n_i(x)$ on a full set $G_f$. If forward sequential
Bowen-mass comparability holds, then $\widetilde\mu$ is blockwise Bowen
sequential weak Gibbs for $(F,\varphi)$ with the same sequential pressure $P$.
For $\widetilde\mu$-a.e. $\hat x$, the prescribed sequence
$n_i(\pi\hat x)$ remains admissible upstairs.
\end{theorem}

\begin{proof}
Fix $\hat x\in G_F$, put $x=\pi\hat x$, choose a Gibbs time
$n_i=n_i(x)$, and let $0\le j<n_i$. Put $m=n_i-j$. The cohomological identity
gives
\begin{equation}
 S_m\varphi(F^j\hat x)
 =S_m\psi(f^jx)+u(F^j\hat x)-u(F^{n_i}\hat x),
 \label{eq:seq-birkhoff-coboundary}
\end{equation}
so
\begin{equation}
 \left|S_m\varphi(F^j\hat x)-S_m\psi(f^jx)\right|
 \le2\|u\|_\infty.
 \label{eq:seq-birkhoff-error}
\end{equation}
By Definition~\ref{def:bowen-sequential-weak}, the same sequence $n_i(x)$
works simultaneously at the reparametrized radii $\eps_-$ and $\eps_+$.
Combining those estimates with \eqref{eq:seq-forward-comparison} gives a
Gibbs constant upstairs bounded by
\[
 \widetilde K_i(\hat x,\eps)
 \le
 \exp\bigl(a_i(\hat x,\eps)+2\|u\|_\infty\bigr)
 \max\{K_i(x,\eps_-),K_i(x,\eps_+)\}.
\]
The same number $P$ appears in both exponential weights because the
cohomological correction contributes only the bounded endpoint term in
\eqref{eq:seq-birkhoff-coboundary}. Therefore
\[
 \frac1{n_i(x)}\log \widetilde K_i(\hat x,\eps)\longrightarrow0.
\]
The estimate is uniform over all suffixes $0\le j<n_i(x)$, so the prescribed
sequence remains admissible upstairs.
\end{proof}

\begin{theorem}[Backward blockwise sequential weak Gibbs transfer]
\label{thm:sequential-weak-backward}
Assume
\[
 \varphi=\psi\circ\pi+u-u\circ F,
\]
and let $\widetilde\mu\leftrightarrow\mu$. Suppose $\widetilde\mu$ is
blockwise Bowen sequential weak Gibbs for $(F,\varphi)$ with sequential
pressure $P$ on a full set $G_F$, and let $\ell:G_f\to G_F$ be a prescribed
lift as in Definition~\ref{def:sequential-backward-comparison}. If backward
sequential Bowen-mass comparability along $\ell$ holds, then $\mu$ is
blockwise Bowen sequential weak Gibbs for $(f,\psi)$ with the same sequential
pressure $P$, along the admissible sequence
\[
 n_i(x)=n_i(\ell(x)).
\]
\end{theorem}

\begin{proof}
Fix $x\in G_f$, write $\hat x=\ell(x)$, and let
$n_i=n_i(\hat x)$. For $0\le j<n_i$ put $m=n_i-j$. The identity
\eqref{eq:seq-birkhoff-coboundary}, applied to $\hat x$, gives the same bounded
cohomological error \eqref{eq:seq-birkhoff-error}. Combining the sequential
weak Gibbs estimates for $\widetilde\mu$ at the radii $\eps_-$ and $\eps_+$
with \eqref{eq:seq-backward-comparison} yields a quotient Gibbs constant
bounded by
\[
 K_i^{\mathrm{base}}(x,\eps)
 \le
 \exp\bigl(b_i(x,\eps)+2\|u\|_\infty\bigr)
 \max\{\widetilde K_i(\hat x,\eps_-),\widetilde K_i(\hat x,\eps_+)\}.
\]
Dividing the logarithm by $n_i(x)$ and using the two subexponential
hypotheses gives the desired blockwise sequential weak Gibbs estimate.
\end{proof}

\begin{corollary}[Strong Bowen sequential Gibbs transfer]
\label{cor:sequential-gibbs-transfer}
Under the hypotheses of Theorem~\ref{thm:sequential-weak-forward}, assume
that the quotient measure is Bowen sequential Gibbs in the sense of
Definition~\ref{def:bowen-sequential-gibbs} and that forward sequential
Bowen-mass comparability is uniformly bounded. Then the lifted measure is
Bowen sequential Gibbs with the same sequential pressure, and the prescribed
admissible Gibbs-time sequence used in the forward comparison remains
admissible upstairs.

Conversely, under the hypotheses of Theorem~\ref{thm:sequential-weak-backward},
if the extension measure is Bowen sequential Gibbs and the backward
comparison is uniformly bounded, then the quotient measure is Bowen
sequential Gibbs along the sequence induced by the prescribed lift.
\end{corollary}

\begin{proof}
In the forward direction, for each fixed radius the proof of
Theorem~\ref{thm:sequential-weak-forward} gives a Gibbs constant bounded by
\[
 e^{A(\eps)+2\|u\|_\infty}
 \max\{K(\eps_-),K(\eps_+)\},
\]
where $A(\eps)$ is a uniform bound for the logarithmic mass error. This
constant is independent of the Gibbs index and the suffix. The backward
statement follows identically from
Theorem~\ref{thm:sequential-weak-backward}.
\end{proof}

\begin{remark}[Why subexponential and bounded comparison differ]
An $o(n_i)$ logarithmic mass error is sufficient for the blockwise sequential
weak Gibbs transfer, but it does not in general preserve a fixed Gibbs
constant. Uniformly bounded mass distortion is therefore the natural
hypothesis for preserving the strong Bowen sequential Gibbs property.
\end{remark}

\begin{remark}[Relation with factor sequential Gibbs measures]
\label{rem:factor-sequential-gibbs}
Ferreira and Oliveira \cite{FerreiraOliveira} prove, for regular one-block
factor maps between one-sided symbolic systems, that the image of a sequential
Gibbs measure is sequential Gibbs and that the maximal Gibbs-time sequence is
preserved in their cylinder framework. Their result extends the classical
factor-Gibbs problem studied, among others, in
\cite{KemptonPollicott,ChazottesUgalde,KemptonSFT,Piraino}. Their result is a projection theorem
and constructs a potential on the symbolic factor. The mechanism above is
complementary: for a fiber-contracting extension it gives a lifting theorem
once the potential has been canonically reduced to the quotient and the
required Bowen-mass comparison is available.

Our Bowen formulation should not be identified with the cylinder formulation
without an additional symbolic comparison lemma. In particular, for two-sided
shifts a small Bowen ball is a bilateral cylinder whose time window is enlarged
by a radius-dependent but fixed number of coordinates. Thus the relation with
Ferreira--Oliveira is exact only in symbolic settings where this finite-window
change has uniformly bounded measure distortion along the relevant Gibbs
blocks.
\end{remark}

\begin{lemma}[Cylinder-to-Bowen bridge under bounded finite-window distortion]
\label{lem:cylinder-bowen-bridge}
Let $(\Sigma,\sigma)$ be a symbolic system endowed with a standard symbolic
metric, and let $\mu$ be a probability measure. Suppose that for every
sufficiently small Bowen radius $\eps>0$ there is a constant
$C(\eps)\ge1$ such that, for every relevant orbit block, every $n\ge1$, and
every $0\le j<n$,
\begin{equation}
 C(\eps)^{-1}
 \mu([x_j\ldots x_{n-1}])
 \le
 \mu\bigl(B_\sigma(\sigma^jx,\eps,n-j)\bigr)
 \le
 C(\eps)
 \mu([x_j\ldots x_{n-1}]).
 \label{eq:cylinder-bowen-suffix-comparison}
\end{equation}
Here $[x_j\ldots x_{n-1}]$ denotes the suffix cylinder used in the symbolic
sequential Gibbs definition. In a two-sided shift, condition
\eqref{eq:cylinder-bowen-suffix-comparison} is, for instance, implied when a
Bowen ball is the corresponding radius-dependent bilateral enlargement of
that suffix cylinder and the measure of this enlarged cylinder is uniformly
comparable, by a factor depending only on $\eps$, with the measure of the
central suffix cylinder.

If the cylinder estimates defining sequential Gibbs hold along a prescribed
admissible sequence $n_i(x)$, then the same sequence is admissible for the
Bowen sequential Gibbs definition, with Gibbs constants multiplied by at most
$C(\eps)$.
\end{lemma}

\begin{proof}
Fix a Gibbs block $n_i(x)$ and a suffix $0\le j<n_i(x)$. Apply the symbolic
cylinder Gibbs estimate to $[x_j\ldots x_{n_i(x)-1}]$ and then use
\eqref{eq:cylinder-bowen-suffix-comparison} with $n=n_i(x)$. The resulting
upper and lower Bowen-ball estimates acquire at most the multiplicative factor
$C(\eps)$. Since this factor is independent of the Gibbs index $i$ and of the
suffix $j$, the same prescribed sequence remains admissible in the Bowen
formulation.
\end{proof}

\section{Explicit Examples and Applications}
\label{sec:examples}

The preceding results are abstract by design. This section verifies the
hypotheses and computes the reduction explicitly in concrete models. The
first example is the basic test case for the whole framework: the potential
depends genuinely on the contracting coordinate, yet the transfer function
and quotient potential can be written in closed form. The second example
shows how the same computation interacts with Gibbs measures. We then
indicate which conclusions can be drawn for solenoids and partially
hyperbolic attractors only conditionally.

\subsection{Affine contracting skew-products: an explicit reduction}

Let $M=\mathbb S^1=\mathbb R/\mathbb Z$ and let
\[
f(x)=dx\pmod 1,\qquad d\ge2.
\]
Fix $R>0$, $|a|<1$, and a Lipschitz function
$\rho:\mathbb S^1\to\mathbb R$ satisfying
\[
\|\rho\|_\infty\le (1-|a|)R.
\]
Then
\[
F:\mathbb S^1\times[-R,R]\longrightarrow
\mathbb S^1\times[-R,R],
\qquad
F(x,y)=\bigl(f(x),ay+\rho(x)\bigr)
\]
is a well-defined continuous skew-product. We take
\[
\pi(x,y)=x.
\]
The fibers are vertical intervals and
\[
|y_n-y_n'|=|a|^n|y-y'|
\]
whenever the initial points have the same base coordinate. Thus Assumption
\ref{ass:selector} holds with the continuous section
\[
s(x)=(x,0),
\qquad
r(x,y)=(x,0).
\]

The next proposition computes the reduction for a class of potentials that
is simple enough to be explicit but is not constant along the contracting
fibers.

\begin{proposition}[Explicit affine reduction]
\label{prop:affine-explicit}
Let
\[
F(x,y)=\bigl(f(x),ay+\rho(x)\bigr),\qquad |a|<1,
\]
as above, and let
\[
\varphi(x,y)=g(x)+by,
\]
where $g:\mathbb S^1\to\mathbb R$ is continuous and $b\in\mathbb R$.
Then the stable-difference series associated with the selector $s(x)=(x,0)$
converges explicitly and yields
\[
\boxed{u(x,y)=\frac{b}{1-a}\,y}.
\]
If $g$ is H\"older, this is precisely the transfer function furnished by
Theorem~\ref{thm:reduction}, and the reduced potential is
\[
\boxed{\psi(x)=g(x)+\frac{b}{1-a}\rho(x).}
\]
Consequently
\[
\boxed{
g(x)+by
=
g(x)+\frac{b}{1-a}\rho(x)
+\frac{b}{1-a}y
-\frac{b}{1-a}\bigl(ay+\rho(x)\bigr).
}
\]
\end{proposition}

\begin{proof}
For points in the same fiber,
\[
F^j(x,y)=\bigl(f^j(x),y_j\bigr),
\]
where
\[
y_j=a^jy+
\sum_{k=0}^{j-1}a^{j-1-k}\rho(f^kx).
\]
For the selected point $(x,0)$ the corresponding fiber coordinate is
\[
y_j^0=
\sum_{k=0}^{j-1}a^{j-1-k}\rho(f^kx).
\]
Hence
\[
\varphi(F^j(x,y))-\varphi(F^j(x,0))
=b(y_j-y_j^0)=ba^jy.
\]
The stable-difference series therefore gives
\[
u(x,y)=\sum_{j=0}^\infty ba^jy
=\frac{b}{1-a}y.
\]
Finally,
\[
\begin{aligned}
\varphi-u+u\circ F
&=g(x)+by-\frac{b}{1-a}y
+\frac{b}{1-a}(ay+\rho(x))\\
&=g(x)+\frac{b}{1-a}\rho(x),
\end{aligned}
\]
which depends only on $x$.
\end{proof}

\begin{remark}[Why this example matters]
The potential $\varphi(x,y)=g(x)+by$ is not stable-constant unless $b=0$.
Nevertheless, its entire dependence on the contracting coordinate is a
coboundary plus an explicit correction
\[
\frac{b}{1-a}\rho(x)
\]
on the quotient. Thus the reduction is not merely a formal statement that
one may ``assume'' stable constancy: it gives an explicit formula for the
new quotient potential.
\end{remark}

\subsubsection{Consequences for ergodic optimization}

For the affine model,
\[
m(\varphi,F)
=
\min_{\mu\in\M_f(\mathbb S^1)}
\int_{\mathbb S^1}
\left(g+\frac{b}{1-a}\rho\right)\,d\mu.
\]
Thus a minimization problem for the genuinely two-variable observable
$g(x)+by$ is exactly the one-dimensional optimization problem for
\[
\psi=g+\frac{b}{1-a}\rho.
\]
Moreover, minimizing measures correspond bijectively. In particular, if the
quotient potential has a unique minimizing measure $\mu_{\min}$, then the
extension has a unique minimizing measure, namely its unique invariant lift.

For $0\le a<1$, the factor $1/(1-a)$ has the expected memory
interpretation: stronger contraction (smaller $a$) reduces the cumulative
contribution of the fiber coordinate, whereas as $a\to1^-$ that accumulated
contribution becomes large. For negative $a$ the alternating signs change
this monotone interpretation, although the explicit formula remains valid
for every $|a|<1$.

\subsubsection{Consequences for pressure and equilibrium states}

Suppose $g$ and $\rho$ are H\"older. Then
\[
\psi(x)=g(x)+\frac{b}{1-a}\rho(x)
\]
is H\"older. Since $f(x)=dx\pmod1$ is an expanding circle map, the standard
thermodynamic formalism for H\"older potentials gives an equilibrium state
for $\psi$. Theorem~\ref{thm:pressure} then yields
\[
P_F(\varphi)
=
P_f\!\left(g+\frac{b}{1-a}\rho\right),
\]
and the equilibrium state for the original fiber-dependent potential is the
unique invariant lift of the equilibrium state of the quotient potential
whenever uniqueness holds on the base.

Hence the thermodynamic problem for the two-dimensional skew-product has
been reduced exactly to a one-dimensional expanding map.

\subsection{A symbolic skew-product and weak Gibbs transfer}

We next consider a model in which the invariant measure and the Gibbs
correspondence can be seen directly. Let
\[
M=\{0,1\}^{\mathbb Z}
\]
with the left shift $\sigma$. Let $|a|<1$ and let
$\rho:M\to\mathbb R$ be H\"older. On
\[
N=M\times[-R,R]
\]
consider
\[
F(\omega,y)=\bigl(\sigma\omega,ay+\rho(\omega)\bigr),
\]
with $R$ chosen so that $F(N)\subset N$. Again
\[
\pi(\omega,y)=\omega
\]
has the continuous section $s(\omega)=(\omega,0)$ and the fibers are
uniformly contracted.

Let $g:M\to\mathbb R$ be H\"older and, for
\[
\varphi(\omega,y)=g(\omega)+by,
\]
the same computation as in Proposition~\ref{prop:affine-explicit} gives
\[
u(\omega,y)=\frac{b}{1-a}y,
\qquad
\psi(\omega)=g(\omega)+\frac{b}{1-a}\rho(\omega).
\]

There is also an invariant graph. Define
\[
h(\omega)
=
\sum_{k=1}^{\infty}a^{k-1}\rho(\sigma^{-k}\omega).
\]
The series converges uniformly, and
\[
h(\sigma\omega)=ah(\omega)+\rho(\omega).
\]
Therefore
\[
F(\omega,h(\omega))
=
(\sigma\omega,h(\sigma\omega)),
\]
so
\[
\Gamma_h=\{(\omega,h(\omega)):\omega\in M\}
\]
is an invariant graph.

Let $\mu$ be an invariant Gibbs measure for the H\"older potential
$\psi$ on the full shift and let
\[
\widetilde\mu=(\omega\mapsto(\omega,h(\omega)))_*\mu.
\]
Then $\widetilde\mu$ is the unique invariant lift of $\mu$ and is supported
on $\Gamma_h$.

\begin{proposition}[Zero-distortion Gibbs transfer on the invariant graph]
\label{prop:symbolic-gibbs-example}
In the symbolic skew-product above, the pair
$(\widetilde\mu,\mu)$ satisfies both forward and backward Bowen-mass
comparability with zero logarithmic mass error (after a change of Bowen-ball
radii). Consequently, if $\mu$ is Gibbs
for $(\sigma,\psi)$, then $\widetilde\mu$ is weak Gibbs for $(F,\varphi)$.
Conversely, weak Gibbs for $(F,\varphi)$ on the invariant graph implies weak
Gibbs for $(\sigma,\psi)$.
\end{proposition}

\begin{proof}
The graph map
\[
\iota(\omega)=(\omega,h(\omega))
\]
is a homeomorphism from $M$ onto $\Gamma_h$ and conjugates $\sigma$ with
$F|_{\Gamma_h}$. Hence Bowen balls on the graph are carried into Bowen balls
for the shift, with only a change of radius determined by the uniform
continuity of $\iota$ and $\iota^{-1}$. Since
$\widetilde\mu=\iota_*\mu$, no multiplicative measure distortion occurs.
Thus the logarithmic mass errors can be taken equal to zero; only the Bowen-ball radii are reparametrized.
The weak Gibbs transfer follows from Theorems
\ref{thm:weak-gibbs-forward} and \ref{thm:weak-gibbs-backward}.
\end{proof}

\begin{corollary}[Bowen sequential Gibbs on the invariant graph]
\label{cor:symbolic-sequential-gibbs}
In the symbolic skew-product above, if $\mu$ is Bowen sequential Gibbs for
$(\sigma,\psi)$ along an admissible sequence $n_i(\omega)$, then
$\widetilde\mu$ is Bowen sequential Gibbs for $(F,\varphi)$ and the same
prescribed sequence remains admissible on the invariant graph. Conversely,
Bowen sequential Gibbs on the invariant graph transfers to the base along the
sequence induced by the graph inverse.
\end{corollary}

\begin{proof}
The graph conjugacy has zero logarithmic mass distortion, as proved in
Proposition~\ref{prop:symbolic-gibbs-example}; hence the associated forward
and backward sequential Bowen-mass comparisons are uniformly bounded. Apply
Corollary~\ref{cor:sequential-gibbs-transfer}.
\end{proof}

\begin{remark}[Connection with Ferreira--Oliveira in the symbolic example]
If, in addition, a cylinder sequential Gibbs measure in the sense of
Ferreira--Oliveira satisfies the bounded finite-window distortion hypothesis of
Lemma~\ref{lem:cylinder-bowen-bridge}, then its prescribed cylinder Gibbs-time
sequence is also admissible for the Bowen formulation and the preceding
corollary applies. This extra hypothesis is automatic in several classical
Gibbs/quasi-Bernoulli symbolic settings, but it is not a formal consequence of
sequential Gibbs alone.
\end{remark}

\begin{remark}
This example is useful because it verifies the extra hypothesis required
for weak Gibbs transfer rather than assuming it abstractly. Zero relative
entropy alone would not provide this conclusion. Here the invariant graph
identifies the lifted dynamics with the quotient dynamics on the support of
the lifted measure, so the Bowen-ball comparison is exact up to the harmless
change of metric radii.
\end{remark}

\subsection{A direct subcohomological application}

The affine model also provides a concrete illustration of the critical
subcohomological transfer. Suppose the quotient map $f$ belongs to a class
for which property $\mathrm{SC}$ is known, and assume
\[
m\!\left(g+\frac{b}{1-a}\rho,f\right)\ge0.
\]
Then
\[
g+\frac{b}{1-a}\rho\ge v-v\circ f
\]
for some continuous $v$. Proposition~\ref{prop:affine-explicit} gives
\[
g(x)+by
\ge
\left(
\frac{b}{1-a}y+v(x)
\right)
-
\left(
\frac{b}{1-a}(ay+\rho(x))+v(fx)
\right).
\]
Thus the subaction on the original two-dimensional system is explicitly
\[
U(x,y)=\frac{b}{1-a}y+v(x).
\]
This is a particularly transparent instance of the transfer principle:
the difficult subcohomological statement is proved on the quotient, while
the stable coordinate is restored by an explicit coboundary.

If the quotient is a zooming local homeomorphism satisfying the hypotheses
of Mbarki--Santana, the same argument applies to their subcohomological
theorem. The resulting statement is genuinely about the original
fiber-dependent potential, not merely about a potential assumed to be
constant on stable fibers.

\subsection{A zooming/non-uniformly expanding application}\label{subsec:zooming-application}

We now turn the affine model into a genuine application to non-uniformly
expanding quotient dynamics. Let $f:M\to M$ be a continuous zooming local
homeomorphism in a class for which the subcohomological theorem of Mbarki--
Santana applies, and assume that the zooming set is dense in $M$. Let
$\rho,g:M\to\mathbb R$ be H\"older continuous and choose $0<|a|<1$ and
$R>0$ so that
\[
\frac{\|\rho\|_\infty}{1-|a|}\le R.
\]
On $N=M\times[-R,R]$ consider
\[
F(x,y)=\bigl(f(x),ay+\rho(x)\bigr),
\qquad \pi(x,y)=x.
\]
Then $F(N)\subset N$, the fibers are uniformly contracted by $|a|$, and
$s(x)=(x,0)$ is a continuous global section. For the genuinely fiber-dependent
potential
\[
\varphi(x,y)=g(x)+by,
\]
Proposition~\ref{prop:affine-explicit} gives
\[
 u(x,y)=\frac{b}{1-a}y,
 \qquad
 \psi(x)=g(x)+\frac{b}{1-a}\rho(x).
\]
Since $g$ and $\rho$ are H\"older, the reduced potential $\psi$ has the
same H\"older regularity. Therefore every hypothesis of the quotient
subcohomological theorem is checked directly on the reduced potential.

\begin{theorem}[Zooming extension theorem]
\label{thm:zooming-extension}
Assume that the quotient map $f$ and the induced potential $\psi$ satisfy
all hypotheses of the subcohomological theorem of Mbarki--Santana, including
the precise regularity and dynamical assumptions required in that theorem. If
\[
m(\varphi,F)\ge0,
\]
then there exists a continuous function $U:N\to\mathbb R$ such that
\[
\varphi\ge U-U\circ F.
\]
More explicitly, if
\[
\psi=g+\frac{b}{1-a}\rho
\ge v-v\circ f
\]
for a continuous $v:M\to\mathbb R$, then one may take
\[
\boxed{U(x,y)=\frac{b}{1-a}y+v(x).}
\]
If the quotient theorem gives a H\"older subaction $v$, then $U$ has the
corresponding H\"older regularity.
\end{theorem}

\begin{proof}
By the affine reduction,
\[
\varphi=\psi\circ\pi+u-u\circ F,
\qquad u(x,y)=\frac{b}{1-a}y.
\]
The invariant-measure correspondence gives
$m(\psi,f)=m(\varphi,F)\ge0$. The Mbarki--Santana theorem therefore yields
$v$ with $\psi\ge v-v\circ f$. Pulling this inequality back by $\pi$ and
adding the coboundary $u-u\circ F$ gives
\[
\varphi\ge (v\circ\pi+u)-(v\circ\pi+u)\circ F.
\]
Thus $U=v\circ\pi+u$, which is the stated formula.
\end{proof}

\begin{corollary}[Zooming Liv\v{s}ic transfer]
\label{cor:zooming-liv-trans}
In the setting of Theorem~\ref{thm:zooming-extension}, suppose in addition
that
\[
\int_M\psi\,d\mu=0
\quad\text{for every }\mu\in\mathcal M_f(M),
\]
and that the periodic points of $f$ are dense. Assume either that all
hypotheses of the Liv\v{s}ic conclusion of Mbarki--Santana
\cite{MbarkiSantana} are satisfied by the induced potential $\psi$, or, more
abstractly, that the quotient has property $\mathrm{SC}$ for both $\psi$ and
$-\psi$. Then
\[
\varphi=U-U\circ F
\]
for a continuous $U:N\to\mathbb R$. In particular, the fiber-dependent
potential has no additional cohomological obstruction beyond the quotient.
\end{corollary}

\begin{proof}
The invariant-average correspondence implies that all invariant averages of
$\varphi$ vanish exactly when all invariant averages of $\psi$ vanish. Under
the first alternative, the Liv\v{s}ic theorem of Mbarki--Santana gives
$\psi=v-v\circ f$. Under the second alternative,
Corollary~\ref{cor:two-sided-sc-livsic} gives the same
conclusion from the two subcohomological inequalities and density of periodic
points. In either case the reduction formula yields
$U=v\circ\pi+u$.
\end{proof}

\begin{remark}[Why this is a substantive application]
The point is that $\varphi$ is not assumed to be constant along the
contracting fibers. The quotient theorem is applied only after the explicit
cohomological correction has produced $\psi$. Thus the result converts a
subcohomological/Liv\v{s}ic theorem for a non-uniformly expanding quotient
into a theorem for a higher-dimensional fiber extension. The construction is
fully explicit in the fiber coordinate.
\end{remark}

\subsection{What the examples do and do not prove}

The preceding examples verify all hypotheses of the reduction theorem and,
in the symbolic case, the additional hypotheses required for weak Gibbs
transfer. They should be contrasted with the following important classes.

\subsubsection{Smale--Williams solenoids}

The classical Smale--Williams solenoid has contracting stable fibers and an
expanding quotient model \cite{Williams}. The present framework applies to a
chosen quotient representation only if its factor projection admits a
continuous global section satisfying Assumption~\ref{ass:selector}. This
condition is topological and is not implied merely by hyperbolicity.
Accordingly, the solenoid is a natural target for the theory, but we do not
claim a general solenoid theorem without verifying that additional
hypothesis.

\subsubsection{Partially hyperbolic attractors}

Several partially hyperbolic attractors admit quotient descriptions and
have well-developed thermodynamic theories
\cite{FisherOliveira,RamosSiqueira,RiosSiqueira,
ClimenhagaPesinZelerowicz,ParmenterPollicott}. Whenever a concrete quotient
projection has a continuous section and uniform contraction in the fibers,
the present reduction can be combined with those quotient results. In this
sense the framework is complementary to approaches that start by assuming
the potential is already constant along stable leaves: here stable
constancy is obtained cohomologically.

The important point is that this subsection states a conditional
application, not a new theorem about all partially hyperbolic attractors.
Removing the global-section requirement, or replacing it by intrinsic
holonomy assumptions, is a natural direction for future work.

\subsubsection{Zooming and non-uniformly expanding quotients}

The preceding affine example can also be used with a quotient map from a
class for which subcohomology or Liv\v{s}ic theory is available, provided
the regularity of the induced potential is verified. The division of labor
is then
\[
\boxed{
\begin{array}{c}
\text{fiber contraction}\\
\Downarrow\\
\text{cohomological reduction}\\
\Downarrow\\
\text{quotient subcohomology/Liv\v{s}ic}\\
\Downarrow\\
\text{result on the attractor}
\end{array}}
\]
For zooming quotients, the critical subcohomological step is supplied by
Mbarki--Santana \cite{MbarkiSantana}. For non-uniformly expanding quotients,
one can similarly combine the reduction with the corresponding quotient
thermodynamic theory \cite{OliveiraViana,Pinheiro,VarandasViana,
AlvesRamosSiqueira}.

\section{Relation with Previous Work}\label{sec:comparison}

The reduction mechanism developed here sits at the intersection of several
well-established theories. We emphasize the precise division between the
classical ingredients and the new transfer framework, since the individual
conclusions below are not intended to be claimed as new in isolation.

\subsection{Sinai-type cohomological reduction}

The construction of a cohomologous potential which is constant along a
contracting or stable direction is closely related to the classical Sinai
cohomology lemma \cite{Sinai}. In the present setting the stable reduction is
obtained from the explicit series in Theorem~\ref{thm:reduction}. The novelty
claimed here is instead the systematic use of this reduction for an abstract
fiber-contracting extension, followed by a simultaneous transfer of
subcohomology, Liv\v{s}ic theory, ergodic optimization and thermodynamic
formalism. Thus the paper should be viewed as a reduction-and-transfer
framework rather than as a replacement for Sinai's lemma.
The distinction is substantive. Sinai's classical setting provides a
cohomological normalization along stable directions in hyperbolic systems;
our theorem is formulated instead for an abstract compact factor map with
uniform fiber contraction and a continuous global section. The present paper
does not claim that the convergent series itself is a new cohomological
identity. The contribution developed in this paper is the theorem-level organization of
this reduction together with the measure-theoretic and thermodynamic
consequences that can subsequently be transferred from the quotient. In
particular, no priority claim over the classical stable-direction
cohomology construction is intended.


\subsection{Subcohomology and Liv\v{s}ic theory}

Mbarki and Santana proved subcohomological and Liv\v{s}ic results for
continuous zooming systems, including important non-uniformly expanding
classes \cite{MbarkiSantana}. We use the current arXiv version (v6, revised
15 April 2025). Their quotient results supply the genuinely nontrivial
critical step in our applications. Once the fiber-dependent
potential has been reduced to $\psi$, their result can be lifted through
Theorem~\ref{thm:reduction}. The logical distinction is therefore
\[
\text{zooming subcohomology on }M
\quad\Longrightarrow\quad
\text{subcohomology on the fiber extension }N.
\]
The affine zooming application in Section~\ref{sec:examples} makes this
transfer explicit.
No strengthening of the Mbarki--Santana theorem on the quotient is
claimed here. In particular, their hypotheses are not weakened or replaced.
The step developed here is geometric: after reducing a genuinely fiber-dependent
potential to $\psi$, the quotient theorem can be pulled back and combined
with the explicit coboundary $u-u\circ F$. Thus the contribution is a
transfer result from an already established quotient theory to a class of
fiber-contracting extensions.

\subsection{Ergodic optimization}

The correspondence between invariant measures and the equality of invariant
averages reduce optimization upstairs to optimization downstairs. This is
compatible with the general cohomological viewpoint of ergodic optimization
and calibrated subactions \cite{BouschJenkinson,GaribaldiBook,
GaribaldiLopesThieullen}. The added point here is geometric: the optimizer for
a fiber-dependent potential is obtained by lifting the optimizer of an
explicit quotient potential.

\subsection{Thermodynamic formalism for partially hyperbolic systems}

There is a substantial literature on partially hyperbolic dynamics and on
equilibrium states for partially hyperbolic attractors and horseshoes
\cite{AlvesBonattiViana,CastroNascimento,FisherOliveira,RamosSiqueira,
RiosSiqueira,ClimenhagaPesinZelerowicz}. More recently,
Parmenter and Pollicott constructed equilibrium states for certain partially
hyperbolic attractors by densities, including a setting with subexponential
contraction in a centre-unstable direction \cite{ParmenterPollicott}. Their
construction is complementary to the present approach: rather than building
an equilibrium state directly on the extension, we reduce the potential and
the variational problem to a quotient whenever the global-section and uniform
fiber-contraction hypotheses hold. In particular, our method can be used as a
preprocessing step before applying an available quotient equilibrium theorem.

\subsection{Weak Gibbs measures}

Weak Gibbs measures in non-uniformly hyperbolic dynamics go back, in
particular, to the work of Yuri \cite{Yuri1999,Yuri2000}, where Gibbs bounds are
relaxed by subexponential distortion. The weak Gibbs part of the present paper
is also connected with recent work on weak Gibbs measures for local
homeomorphisms \cite{FerreiraRamos}. That work develops existence, uniqueness
and statistical consequences on quotient-type systems. Here we address the
complementary question of when a weak Gibbs measure on the quotient survives
passage to a fiber-contracting extension. The required
input is not merely zero relative entropy, but quantitative two-sided
Bowen-mass comparison. The invariant-graph example verifies this hypothesis
explicitly.

\subsection{Sequential Gibbs measures and factor maps}

The preservation of Gibbsian structure under symbolic factor maps has a
substantial literature. For full shifts, important results were obtained by
Kempton--Pollicott and by Chazottes--Ugalde
\cite{KemptonPollicott,ChazottesUgalde}; for subshifts of finite type see
Kempton \cite{KemptonSFT}. H\"older regularity of projected Gibbs states under
fiber-wise mixing factors was sharpened by Piraino \cite{Piraino}, while an
almost-additive extension was developed by Yayama \cite{YayamaAlmostAdditive}.

Sequential Gibbs measures were introduced by Ferreira and Oliveira in a
one-sided symbolic setting designed to capture Gibbs estimates available only
at point-dependent times \cite{FerreiraOliveira}. Their theorem gives a
non-uniform counterpart of the preceding factor-Gibbs results: under their
regularity condition on a one-block factor map, the image of a sequential
Gibbs measure is again sequential Gibbs, with the same maximal Gibbs-time
sequence in their cylinder framework. The sequential results of
Section~\ref{sec:sequential-gibbs} address a complementary direction: after
the contracting fiber has been removed cohomologically, sequential Gibbs
information can be lifted from the quotient to the extension whenever
Bowen-ball masses are comparable along the prescribed Gibbs blocks.  Subexponential comparison yields blockwise Bowen sequential weak Gibbs,
while uniformly bounded comparison preserves the strong Bowen sequential
Gibbs property along the prescribed admissible sequence. In symbolic settings,
comparison with the cylinder formulation of Ferreira--Oliveira additionally
uses Lemma~\ref{lem:cylinder-bowen-bridge}.

Thus a fiber-contraction reduction may be composed with a symbolic factor
map: the former is a cohomological lifting mechanism, whereas the latter is
a projection theorem for sequential Gibbs measures.  No claim is made that
the symbolic factor theorem itself follows from fiber contraction.

\subsection{Summary of the contribution}

The contribution of the present paper can therefore be summarized by the
following diagram:
\[
\boxed{
\begin{array}{c}
\text{fiber contraction + continuous section}\\\
\Downarrow\\\
\text{cohomological reduction }\varphi=\psi\circ\pi+u-u\circ F\\\
\Downarrow\\\
\text{apply quotient theory to }\psi\\
\Downarrow\\\
\begin{array}{c}
\text{subcohomology / Liv\v{s}ic}\\
\text{optimization / pressure / equilibrium}\\
\text{weak Gibbs, when Bowen-mass comparison holds}\\
\text{sequential Gibbs, along controlled Gibbs times}
\end{array}
\end{array}}
\]
This is the organizing principle of the paper.

\section{Scope and Further Directions}

The preceding results separate the hypotheses needed for reduction,
regularity, thermodynamic consequences, and weak Gibbs transfer.

\subsection{Relation with neighboring theories}

The reduction theorem is best viewed as a geometric preprocessing step: it
removes the contracting-fiber variation before one invokes a theorem on the
quotient. In ergodic optimization, the identity
\[
 m(\varphi,F)=m(\psi,f)
\]
and the bijection of minimizing measures show that the optimization problem is
completely encoded by the quotient potential. In Liv\v{s}ic theory, automatic
periodic lifting converts the periodic obstruction upstairs into the usual
periodic obstruction downstairs. In thermodynamic formalism, zero relative
entropy explains why the fibers contribute no exponential orbit complexity,
while the bounded coboundary identity
\[
 S_n\varphi-S_n\psi\circ\pi=u-u\circ F^n
\]
accounts for the pressure correspondence. These conclusions use different
parts of the hypotheses: zero relative entropy alone does not imply convergence
of the stable cohomological series, and continuous cohomological reduction
alone does not imply the Bowen-ball comparison required for weak Gibbs
transfer. This separation is essential when applying the framework to concrete
fibered systems.

Several extensions remain natural. One is to replace uniform contraction by
non-uniform or tempered contraction and determine whether the stable series
still defines a continuous transfer. A second is to seek intrinsic geometric
conditions guaranteeing a continuous or H\"older section of the quotient
map. A third is to extend the reduction from scalar observables to vector or
matrix cocycles. Finally, the bounded identity
\[
S_n\varphi-S_n\psi\circ\pi=u-u\circ F^n
\]
suggests that large-deviation statements on the quotient may transfer once
the relevant measure comparison is established. None of these extensions is
used in the proofs above.

\section{Concluding Remarks}

The main structural result of the paper is the cohomological reduction
\[
\boxed{
\varphi=\psi\circ\pi+u-u\circ F.
}
\]
Its significance is that a potential which genuinely varies in the
contracting direction can be replaced, without changing its cohomology
class, by a quotient potential. The subsequent theories are therefore
organized around the quotient rather than proved independently on the
higher-dimensional extension.

The affine skew-product in Section~\ref{sec:examples} makes this mechanism
explicit:
\[
F(x,y)=(f(x),ay+\rho(x)),
\qquad
\varphi(x,y)=g(x)+by
\]
gives
\[
u(x,y)=\frac{b}{1-a}y,
\qquad
\psi(x)=g(x)+\frac{b}{1-a}\rho(x).
\]
Thus the contracting coordinate is not simply discarded; its cumulative
effect is transferred to the quotient potential. This calculation also
shows how the abstract pressure, equilibrium, ergodic-optimization, and
subcohomological statements become concrete statements for an explicit
two-dimensional dynamical system.

The symbolic skew-product example further shows that the weak Gibbs theory
is effective when the additional Bowen-mass hypothesis is geometrically
verified. On an invariant graph the lifted measure is supported on a graph
conjugate to the quotient, so the Bowen-ball masses are compared with no
multiplicative distortion, only with a change of radii. Thus a Gibbs measure
on the quotient lifts to a weak Gibbs measure for the original
fiber-dependent potential.

The resulting transfer principle may be summarized as
\[
\boxed{
\text{cohomological theory of the quotient}
\quad\Longrightarrow\quad
\text{cohomological theory of the attractor},
}
\]
with the corresponding variational and thermodynamic statements transferred
at the same time. The hypotheses are deliberately separated: fiber
contraction gives the stable series and zero relative entropy; additional
regularity is needed for H\"older transfers; quotient subcohomology or
Liv\v{s}ic results are needed for the corresponding cohomological
conclusions; Bowen-mass comparison is needed for weak Gibbs transfer; and
sequential Bowen-mass comparison along the prescribed Gibbs blocks is needed
for the sequential transfer results.

The examples also indicate the principal limitation of the current
framework. A continuous global section is a strong hypothesis and is not
automatic for general solenoidal or partially hyperbolic quotients.
Developing an intrinsic version based on local sections, stable holonomies,
or measurable selectors would substantially enlarge the scope of the
theory. Other natural directions include non-uniform or tempered fiber
contraction, vector and matrix cocycles, and the transfer of large-deviation
principles through suitable quantitative measure comparisons.

\section*{Acknowledgements}
This work was supported by PRAPG\_CAPES-BRASIL.

\begin{thebibliography}{99}

\bibitem{AlvesBonattiViana}
J.~F. Alves, C. Bonatti and M. Viana,
\newblock SRB measures for partially hyperbolic systems whose central
direction is mostly expanding,
\newblock \emph{Invent. Math.} 140 (2000), 351--398.

\bibitem{AlvesRamosSiqueira}
J.~F. Alves, V. Ramos and J. Siqueira,
\newblock Equilibrium stability for non-uniformly hyperbolic systems,
\newblock \emph{Ergodic Theory Dynam. Systems} 39 (2019), 2619--2642.

\bibitem{BouschJenkinson}
T. Bousch and O. Jenkinson,
\newblock Cohomology classes of dynamically non-negative $C^k$ functions,
\newblock \emph{Invent. Math.} 148 (2002), 207--217.

\bibitem{Bowen}
R. Bowen,
\newblock \emph{Equilibrium States and the Ergodic Theory of Anosov
Diffeomorphisms},
\newblock Lecture Notes in Mathematics 470, Springer, 1975.

\bibitem{CastroNascimento}
A. Castro and T. Nascimento,
\newblock Statistical properties of the maximal entropy measure for
partially hyperbolic attractors,
\newblock \emph{Ergodic Theory Dynam. Systems} 37 (2017), 1060--1101.

\bibitem{ClimenhagaPesinZelerowicz}
V. Climenhaga, Y. Pesin and A. Zelerowicz,
\newblock Equilibrium measures for some partially hyperbolic systems,
\newblock \emph{J. Mod. Dyn.} 16 (2020), 155--205.


\bibitem{ChazottesUgalde}
J.-R. Chazottes and E. Ugalde,
\newblock On the preservation of Gibbsianness under symbol amalgamation,
\newblock in \emph{Entropy of Hidden Markov Processes and Connections to
Dynamical Systems}, London Math. Soc. Lecture Note Ser. 385,
Cambridge Univ. Press, 2011, 72--97,
\newblock doi:10.1017/CBO9780511819407.003.

\bibitem{deLaLlaveMarcoMoriyon}
R. de la Llave, J.~M. Marco and R. Moriy\'on,
\newblock Canonical perturbation theory of Anosov systems and regularity
results for the Liv\v{s}ic cohomology equation,
\newblock \emph{Ann. of Math.} (2) 123 (1986), 537--611.

\bibitem{FerreiraRamos}
G. Ferreira and V. Ramos,
\newblock Weak Gibbs measures for local homeomorphisms,
\newblock arXiv:2510.20938, 2025.

\bibitem{FerreiraOliveira}
G. Ferreira and K. Oliveira,
\newblock Sequential Gibbs measures and factor maps,
\newblock \emph{J. Stat. Phys.} 172 (2018), no.~3, 833--853,
\newblock doi:10.1007/s10955-018-2061-6.

\bibitem{FisherOliveira}
T. Fisher and K. Oliveira,
\newblock Equilibrium states for certain partially hyperbolic attractors,
\newblock \emph{Nonlinearity} 33 (2020), 3409--3435.

\bibitem{GaribaldiBook}
E. Garibaldi,
\newblock \emph{Ergodic Optimization in the Expanding Case: Concepts,
Tools and Applications},
\newblock Springer, 2017.

\bibitem{GaribaldiLopesThieullen}
E. Garibaldi, A.~O. Lopes and P. Thieullen,
\newblock On calibrated and separating sub-actions,
\newblock arXiv:math/0612593, 2006.


\bibitem{KemptonPollicott}
T.~M.~W. Kempton and M. Pollicott,
\newblock Factors of Gibbs measures for full shifts,
\newblock in \emph{Entropy of Hidden Markov Processes and Connections to
Dynamical Systems}, London Math. Soc. Lecture Note Ser. 385,
Cambridge Univ. Press, 2011, 246--257,
\newblock doi:10.1017/CBO9780511819407.009.

\bibitem{KemptonSFT}
T.~M.~W. Kempton,
\newblock Factors of Gibbs measures for subshifts of finite type,
\newblock \emph{Bull. London Math. Soc.} 43 (2011), no.~4, 751--764,
\newblock doi:10.1112/blms/bdr010.

\bibitem{LedrappierWalters}
F. Ledrappier and P. Walters,
\newblock A relativised variational principle for continuous transformations,
\newblock \emph{J. London Math. Soc.} (2) 16 (1977), 568--576,
\newblock doi:10.1112/jlms/s2-16.3.568.

\bibitem{Livsic}
A.~N. Liv\v{s}ic,
\newblock Cohomology of dynamical systems,
\newblock \emph{Izv. Akad. Nauk SSSR Ser. Mat.} 36 (1972), 1296--1320.

\bibitem{MbarkiSantana}
L. Mbarki and E. Santana,
\newblock Subcohomology and a Liv\v{s}ic theorem for zooming systems,
\newblock arXiv:2208.13209v6, revised 15 April 2025.

\bibitem{OliveiraViana}
K. Oliveira and M. Viana,
\newblock Thermodynamical formalism for robust classes of potentials and
non-uniformly hyperbolic maps,
\newblock \emph{Ergodic Theory Dynam. Systems} 28 (2008), 501--533.

\bibitem{ParmenterPollicott}
D. Parmenter and M. Pollicott,
\newblock Constructing equilibrium states for some partially hyperbolic
attractors via densities,
\newblock \emph{Ergodic Theory Dynam. Systems}, First View (2026), 1--16,
\newblock published online 17 April 2026, doi:10.1017/etds.2026.10299.


\bibitem{Piraino}
M. Piraino,
\newblock Projections of Gibbs states for H\"older potentials,
\newblock \emph{J. Stat. Phys.} 170 (2018), 952--961,
\newblock doi:10.1007/s10955-018-1967-3.

\bibitem{Pinheiro}
V. Pinheiro,
\newblock Expanding measures,
\newblock \emph{Ann. Inst. H. Poincar\'e Anal. Non Lin\'eaire} 28 (2011),
889--939.

\bibitem{RamosSiqueira}
V. Ramos and J. Siqueira,
\newblock On equilibrium states for partially hyperbolic horseshoes:
uniqueness and statistical properties,
\newblock \emph{Bull. Braz. Math. Soc.} 48 (2017), 347--375.

\bibitem{RiosSiqueira}
I. Rios and J. Siqueira,
\newblock On equilibrium states for partially hyperbolic horseshoes,
\newblock \emph{Ergodic Theory Dynam. Systems} 38 (2018), 301--335.

\bibitem{Sinai}
Ya.~G. Sinai,
\newblock Gibbs measures in ergodic theory,
\newblock \emph{Russian Math. Surveys} 27 (1972), no.~4, 21--69,
\newblock doi:10.1070/RM1972v027n04ABEH001383.

\bibitem{Ruelle}
D. Ruelle,
\newblock \emph{Thermodynamic Formalism},
\newblock Addison--Wesley, Reading, MA, 1978.

\bibitem{VarandasViana}
P. Varandas and M. Viana,
\newblock Existence, uniqueness and stability of equilibrium states for
non-uniformly expanding maps,
\newblock \emph{Ann. Inst. H. Poincar\'e Anal. Non Lin\'eaire} 27 (2010),
555--593.

\bibitem{Walters}
P. Walters,
\newblock \emph{An Introduction to Ergodic Theory},
\newblock Graduate Texts in Mathematics 79, Springer, 1982.

\bibitem{Williams}
R.~F. Williams,
\newblock Expanding attractors,
\newblock \emph{Publ. Math. IH\'ES} 43 (1974), 169--203.


\bibitem{YayamaAlmostAdditive}
Y. Yayama,
\newblock On factors of Gibbs measures for almost additive potentials,
\newblock \emph{Ergodic Theory Dynam. Systems} 36 (2016), no.~1, 276--309,
\newblock doi:10.1017/etds.2014.50.

\bibitem{Yuri1999}
M. Yuri,
\newblock Thermodynamic formalism for certain nonhyperbolic maps,
\newblock \emph{Ergodic Theory Dynam. Systems} 19 (1999), no.~5, 1365--1378,
\newblock doi:10.1017/S0143385799141695.

\bibitem{Yuri2000}
M. Yuri,
\newblock Weak Gibbs measures for certain non-hyperbolic systems,
\newblock \emph{Ergodic Theory Dynam. Systems} 20 (2000), no.~5, 1495--1518,
\newblock doi:10.1017/S014338570000081X.

\end{thebibliography}

\end{document}
